\chapter{Master concordance: the seven-face dictionary across three volumes}
\label{ch:master-concordance}

\providecommand{\Tr}{\operatorname{Tr}}

\begin{abstract}
The common datum in the seven-face dictionary is the genus-zero
binary collision residue
\[
 r_\cA(z)
 = \operatorname{Res}^{\mathrm{coll}}_{0,2}(\Theta_\cA),
 \qquad
 \Theta_\cA = D_\cA-d_0.
\]
It is the binary component of the twisting morphism
$\theta_\cA\colon B(\cA)\to\cA$.  Bar--cobar inversion recovers
$\cA$; Verdier duality produces $\cA^!$; chiral Hochschild cochains
produce $Z^{\mathrm{der}}_{\mathrm{ch}}(\cA)$.  The tables below record the
normalizations under which the same collision residue appears as
twisting morphism, line-operator twist, Poisson-vertex $r$-matrix,
commuting Hamiltonian seed, Yangian classical limit, Sklyanin tensor,
and Gaudin generator.
\end{abstract}

\section{Five objects and their comparison morphisms}
\label{sec:master-concordance-five-objects}

\begin{convention}[Five typed objects; \ClaimStatusDefinitional]
\label{conv:master-five-object-firewall}
Let $\cA=T(V)/(R)$ be an augmented quadratic chiral algebra on a
curve.  The five objects are
\[
\begin{array}{rcl}
\cA && \text{the chiral algebra;}\\[2pt]
B(\cA)=T^c(s^{-1}\bar\cA) && \text{the ordered bar coalgebra,
with }\bar\cA=\fib(\epsilon);\\[2pt]
\cA^{\mathrm i}=C(s^{-1}V,s^{-2}R) && \text{the quadratic
coalgebra of the chosen presentation;}\\[2pt]
\cA^!=\mathbb D_{\Ran}(\cA^{\mathrm i}) && \text{the strict
Verdier--Koszul dual algebra in the finite-type dualizable ambient;}\\[2pt]
Z^{\mathrm{der}}_{\mathrm{ch}}(\cA)
 = C^\bullet_{\mathrm{ch}}(\cA,\cA) && \text{the chart cochain
centre.}
\end{array}
\]
They are joined by the typed morphisms
\[
\tau_{\mathrm i}\colon\cA^{\mathrm i}\longrightarrow\cA,
\qquad
q_\cA\colon\cA^{\mathrm i}\longrightarrow B(\cA),
\qquad
\epsilon_\cA\colon\Omega B(\cA)\longrightarrow\cA,
\]
\[
\mathbb D_{\Ran}(q_\cA)\colon
\mathbb D_{\Ran}B(\cA)\longrightarrow\cA^!,
\qquad
\mathfrak Z(\cA):=C^\bullet_{\mathrm{ch}}(\cA,\cA).
\]
The package
$H_{\mathrm{CL}}(\cA,\cA^{\mathrm i},\tau_{\mathrm i})$ makes
$q_\cA$ a quasi-isomorphism;
$H_1$ makes $\epsilon_\cA$ a quasi-isomorphism; and $H_{\mathrm{VD}}$
defines the displayed Verdier dual morphism in the completed ambient.
\end{convention}

\begin{theorem}[The five objects and their ambient categories;
\ClaimStatusProvedHere]
\label{thm:typed-verdier-koszul-firewall}
\index{Verdier--Koszul duality!typed objects}
\textup{[}Type signature:
$(\mathrm{Open},\mathrm{five\ typed\ objects},
\mathrm{levels}=1\leftrightarrow2\to3,
\mathrm{hypothesis\ package}=\mathrm{augmented}
\ +\ \mathrm{complete}
\ +\ \mathrm{Verdier\ dualizable}
\ +\ \mathrm{chiral\ Hochschild\ defined})$\textup{].}
The five objects of Convention~\ref{conv:master-five-object-firewall}
belong to five different ambient categories:
\[
\begin{array}{c|c}
\text{object} & \text{ambient}\\ \hline
\cA &
\operatorname{Alg}^{\mathrm{aug,comp}}_{\chirAss}(D_X)\\[2pt]
B_X^{\mathrm{ord}}(\cA) &
\operatorname{Coalg}^{\mathrm{conil,comp}}_{\chirAss^!}(D_X)\\[2pt]
\cA^{\mathrm i} &
\operatorname{Coalg}^{\mathrm{quad}}_{\chirAss^!}(D_X)\\[2pt]
\cA^! &
\operatorname{Alg}^{\mathrm{post}\text{-}\mathbb D_{\Ran},
\mathrm{comp}}_{\chirAss}(D_X)\\[2pt]
Z^{\mathrm{der}}_{\mathrm{ch}}(\cA) &
\operatorname{Ch}\bigl(C^\bullet_{\mathrm{ch}}(\cA,\cA)\bigr).
\end{array}
\]
The typed morphisms are
\[
B_X^{\mathrm{ord}}(\cA)\to B_X^\Sigma(\cA)
\quad(\text{equivariant descent}),\qquad
\Omega_XB_X^{\mathrm{ord}}(\cA)\to\cA
\quad(\text{bar--cobar counit}),
\]
\[
\mathbb D_{\Ran}B_X^{\mathrm{ord}}(\cA)\rightsquigarrow\cA^!
\quad(\text{Verdier--Koszul construction}),\qquad
C^\bullet_{\mathrm{ch}}(\cA,\cA)\to
Z^{\mathrm{der}}_{\mathrm{ch}}(\cA)
\quad(\text{definition of the derived chiral centre}).
\]
The quadratic inclusion is
\(q_\cA\colon\cA^{\mathrm i}\to B_X^{\mathrm{ord}}(\cA)\).
The counit has target \(\cA\), the Verdier morphism has target
\(\cA^!\), and the open--closed morphism of
Convention~\ref{conv:master-oca} compares physical bulk observables
with \(Z^{\mathrm{der}}_{\mathrm{ch}}(\cA)\).
\end{theorem}

\begin{proof}
The ambient categories have different variance and structure:
augmented algebras, conilpotent coalgebras, quadratic coalgebras,
post-Verdier completed algebras, and cochain complexes.  The bar--cobar
counit is the adjunction counit in the algebra/coalgebra pair and has
target \(\cA\).  Verdier--Koszul duality applies
\(\mathbb D_{\Ran}\) to the bar-side comparison and lands in the dual
algebra.  The derived chiral centre is the Hochschild cochain
construction in the bimodule category.  These constructions give the
displayed sources and targets.
\end{proof}

\begin{proposition}[Terminal ordered-to-symmetric survival diagram;
\ClaimStatusConditional]
\label{prop:master-ordered-symmetric-survival-diagram}
\index{ordered-to-symmetric survival diagram}
\index{averaging map!surviving invariants}
\textup{[}Type signature:
$(\mathrm{Open},\mathrm{ordered/symmetric\ bar},
\mathrm{level}=2\to5,
\mathrm{hypothesis\ package}=\mathrm{equivariant\ descent}
\ +\ \mathrm{strict\ ML\ completion}
\ +\ \mathrm{proved\ trace\ map})$\textup{].}
Let \(\cA\) be an augmented chiral algebra satisfying the stated
completion and equivariance hypotheses, and let
\[
\mathrm{av}\colon B^{\mathrm{ord}}(\cA)\longrightarrow B^\Sigma(\cA)
\]
be the ordered-to-symmetric averaging map.  The terminal information
flow is the commutative diagram
\[
\begin{array}{ccccc}
B^{\mathrm{ord}}(\cA)
& \xrightarrow{\ \mathrm{av}\ }
& B^\Sigma(\cA)
& \xrightarrow{\ \Tr_{\mathrm{mod},\log}\ }
& \{\kappa,S_3,S_4,\ldots\}_{\mathrm{scalar}} \\[4pt]
\Big\downarrow{\scriptstyle \pi_{\mathrm{ord}}}
&&
\Big\downarrow{\scriptstyle \pi_{\Sigma}}
&&
\Big\downarrow{\scriptstyle \mathrm{id}} \\[4pt]
\{r_\cA,\beta_{\mathrm{br}},\Phi_{\mathrm{assoc}},
\Theta_\cA^{\ge3}\}_{\mathrm{ordered}}
& \xrightarrow{\ \mathrm{coinv}\ }
& \{\kappa(\cA),S_3^\Sigma,S_4^\Sigma,\ldots\}
& \xrightarrow{\ \mathrm{trace}\ }
& \{\kappa,S_3,S_4,\ldots\}_{\mathrm{scalar}} .
\end{array}
\]
The diagram asserts the displayed arrows.  Averaging carries the
coinvariant scalar classes through \(\pi_\Sigma\), and the modular trace
carries them to the scalar target.  The ordered \(r\)-matrix, braid and
associator, non-symmetric line operators, Yetter--Drinfeld--Schauenburg
brackets, and the physical and CY realizations remain objects of their
respective source theories.  Their transport is specified by a morphism
with displayed source, target, and hypotheses.
\end{proposition}

\begin{convention}[Primitive structural datum]
\label{conv:master-primitive-datum}
The Open quadrant begins with the pointed factorization datum
\[
 \mathfrak D_{\mathrm{op}}
 =\bigl(X,D,\tau;\mathsf C,b,\varepsilon_b\bigr).
\]
Here $(X,D,\tau)$ is a smooth projective curve with reduced marked
divisor and tangential structure, $\mathsf C$ is a stable presentable
factorization category, $b=(b_U)$ is a factorization compact generator,
and
\[
 \varepsilon_b\colon
 A_b=\RHom_{\mathsf C}(b,b)\longrightarrow\mathbf1_{\mathrm{fact}}
\]
is a morphism of unital factorization $E_1$-algebras.  The six levels
are formed from this datum as follows:
\[
\begin{array}{c|l}
0 & \cC^{\mathrm{op}}=(\mathsf C^\omega)^{\mathrm{op}},\\
1 & (A_b,\varepsilon_b),\\
2 & B_{\varepsilon_b}(A_b)=T^c(s^{-1}\bar A_b),\quad
    \theta_b\colon B_{\varepsilon_b}(A_b)\to A_b,\\
3 & Z_{\mathrm{ch}}^{\mathrm{chart}}(A_b)
    =C^\bullet_{\mathrm{ch}}(A_b,A_b),\\
4 & \mathsf{Line}^{\mathrm{int}}(\mathsf C),\quad
    (\mathrm{RMod}^{\mathrm{fact}}_{A_b^e},A_b),\quad
    \mathsf{Line}^{!}(A_b)=\mathrm{Perf}(A_b^!),\\
5 & \mathfrak F,\quad
    \operatorname{Obs}^{\mathrm{def}}_g(A_b),\quad
    \mathfrak O^K_{g,n}(A_b),\quad
    \operatorname{obs}^{\mathrm{Hdg}}_{g,n}(A_b),\quad
    F_{g,n}(A_b;\alpha).
\end{array}
\]
The level-$3$ physical interpretation is carried by the open--closed
morphism of Convention~\ref{conv:master-oca}.  At level~$5$ the
category-valued modular functor $\mathfrak F$ and its realization are
additional data; its scalar projection produces the displayed Hodge
classes and numerical insertions.
\end{convention}

\begin{convention}[Type-signature discipline; KSDual; vertical
equivalences]
\label{conv:master-type-signature}
Every theorem statement carries a four-component type signature
\[
(\textit{quadrant},\;\textit{presentation},\;\textit{level},\;
\textit{hypothesis package}),
\]
where \emph{quadrant} ranges over
$(\mathrm{Open}\cup\mathrm{CY})\times(\mathrm{Cat}\cup\mathrm{Chain})$,
\emph{presentation} names the chain-level, $(\infty,1)$-categorical,
factorisation, or scalar realization,
\emph{level} names the Beilinson-tower height
$0,1,2,3,4,5$ (or a cross-level pair $i\leftrightarrow j$, $i\to j$),
and \emph{hypothesis package} lists the conditions on the chart
algebra, chosen generator, and ambient category.  A comparison is
written in the form
\[
 \mathcal X\xrightarrow[H]{\ f\ }\mathcal Y
 \qquad(\textit{status};\ \textit{ambient}),
\]
so that its source, target, morphism, assumptions, status, and ambient
category appear together.

The Verdier--Koszul reflection
$\sigma(A)=A^!$ acts on the maximal $\infty$-groupoid of the
finite-type curved, weight-completed ambient.  Its homotopy-fixed locus
over the uncurved Koszul charts is
$\KSDual(X,D,\tau)$ of Definition~\ref{def:mr-ksdual}.  A point is a
triple $(A_b,\varphi_b,H_b)$ with
\[
 \varphi_b\colon A_b\xrightarrow{\sim}A_b^!,
 \qquad
 H_b\colon\sigma(\varphi_b)\circ\varphi_b\simeq\eta_{A_b},
\]
together with the higher $\mathbb Z/2$ coherence.  Theorem~A supplies
the separate bar--cobar equivalence under $H_1$; Theorem~H carries its
family-indexed deformation-retract hypotheses.

The six Open--CY comparison morphisms are the maps
$\Xi_{V0}$, $\chi_{V1}$, $B_C(\chi_{V1})$,
$\mathrm{OCA}_{\mathcal T}$, $\rho_{V4}$, and the level-$5$ modular
and scalar comparisons of
Theorem~\ref{thm:mr-vertical-dictionary}.  Their hypothesis packages
$H_{V0},\ldots,H_{V5}$ state the precise domains on which these maps
are equivalences.
\end{convention}

\begin{convention}[Reconstruction and formation maps on the Beilinson tower]
\label{conv:master-reconstruction-tower}
The Open Beilinson fork consists of the following typed maps:
\begin{center}
\renewcommand{\arraystretch}{1.35}
\begin{tabular}{|p{1.7cm}|p{4.5cm}|p{7.7cm}|}
\hline
\textbf{Levels} & \textbf{Map} & \textbf{Hypothesis package and status} \\
\hline\hline
$0\leftrightarrow1$ &
Factorization Morita &
$U\mapsto\mathsf C(U)$ is stable presentable; its structure functors
are exact, colimit preserving, and compact preserving; the compatible
family $b_U$ compactly generates every chart and satisfies the
factorisation K\"unneth equivalence.  Under $H_0$ this is an
equivalence. \\
\hline
$1\leftrightarrow2$ &
Associative chiral bar--cobar, Theorem~A &
The reduced associative operad lies in the pro-nilpotent
Francis--Gaitsgory Ran ambient.  The package
$H_1=H_{\mathrm{fact}}\cup H_{\mathrm{conv}}$ places the augmented
factorization algebras and divided-power conilpotent factorization
coalgebras in complete chain models.  The unit and counit are
equivalences on this domain.\\
\hline
$1\to3$ &
Chiral cochain-center formation &
$A_b\mapsto Z_{\mathrm{ch}}^{\mathrm{chart}}(A_b)$ is functorial.
Under $H_Z$, the map
$\chi_b\colon Z_{\mathrm{ch}}^{\mathrm{chart}}(A_b)\to
Z_{\mathrm{ch}}^{\mathrm{int}}(\mathsf C)$ is an equivalence. \\
\hline
$1\to4$ &
Intrinsic and Koszul-dual operator formation &
Factorization Eilenberg--Watts gives
$\mathsf{Line}^{\mathrm{int}}(\mathsf C)\simeq
\mathrm{RMod}^{\mathrm{fact}}_{A_b^e}$ under $H_0$.
The packages $H_{\mathrm{CL}}$, $H_{\mathrm{VD}}$, and
$H_{\mathrm{line}}$ form and realize
$\mathsf{Line}^{!}(A_b)=\mathrm{Perf}(A_b^!)$. \\
\hline
$5\to4$ &
Genus-zero restriction and realization &
Brochier--Woike genus-zero restriction applies under $H_5$.
The package $H_{\mathrm{real}}$ specifies an equivalence from the
resulting balanced braided algebra to
$\mathsf{Line}^{\mathrm{int}}(\mathsf C)$. \\
\hline
$5\to\mathrm{Scal}$ &
Hodge and numerical projection &
Under~$H_D^K$,
$\mathfrak O^K_{g,n}(A_b)=\kappa(A_b)\lambda_{-1}(\mathbb E_g)$ and
$\operatorname{ch}_g(\mathfrak O^K_{g,n})
=(-1)^g\kappa(A_b)\lambda_g$.  Under~$H_D^{\mathrm{graph}}$, an
insertion
$\alpha\in H^{2(2g-3+n)}(\overline{\mathcal M}_{g,n})$ gives
$F_{g,n}=F^{\mathrm{sc}}_{g,n}+\delta F^{\mathrm{cross}}_{g,n}$. \\
\hline
\end{tabular}
\end{center}
The first two rows are reconstruction equivalences.  The remaining
rows form the center, operator, modular, and scalar branches.  Their
joint statement is Theorem~\ref{thm:mr-master}; its equivariant
restriction to the homotopy-fixed locus is
Theorem~\ref{thm:mr-ksdual}.
\end{convention}

\begin{convention}[Open--closed comparison; OCA]
\label{conv:master-oca}
The derived chiral centre $Z^{\mathrm{der}}_{\mathrm{ch}}(\cA)
= C^\bullet_{\mathrm{ch}}(\cA,\cA)$ is the chart cochain center of
$\cA=A_b$.  For a holomorphic--topological theory $T$ with boundary
chart algebra $\cA$, an open--closed morphism has type
\[
\mathrm{OCA}_T\colon
\mathcal{O}^{\mathrm{phys}}_{\mathrm{bulk}}(T)
\longrightarrow
Z^{\mathrm{der}}_{\mathrm{ch}}(\cA),
\]
together with compatibility with the
$\mathsf{SC}^{\mathrm{ch,top}}$-action on the pair
$(\mathcal{O}^{\mathrm{phys}}_{\mathrm{bulk}}(T),\cA)$; see Vol~II,
\texttt{thm:sc-chtop-deligne}.  On a sector where
$\mathrm{OCA}_T$ is a quasi-isomorphism, it identifies the physical
bulk observables with the chart cochain center.

The hypothesis package $H_{\mathrm{OCA}}$ consists of the
Costello--Gwilliam factorization model on the
holomorphic--topological locus, the local formal-Darboux comparison of
\texttt{thqg\_open\_closed\_realization.tex}, global descent, and the
displayed $\mathsf{SC}^{\mathrm{ch,top}}$ compatibility.  The physical
bulk identification carries status $\ClaimStatusConditional$ on the
domain specified by $H_{\mathrm{OCA}}$.
\end{convention}

\begin{proposition}[Physical-bulk comparison through the open--closed map;
\ClaimStatusConditional]
\label{prop:master-oca-holography-gate}
\index{physical bulk observables!open--closed comparison}
\index{holography!open--closed quasi-isomorphism}
\textup{[}Type signature:
$(\mathrm{Open}\leftrightarrow\mathrm{HT\ bulk},
\mathrm{chain\ comparison},\mathrm{level}=3,
\mathrm{hypothesis\ package}=\mathrm{OCA}_T
\ \mathrm{quasi\text{-}isomorphism}
\ +\ \mathsf{SC}^{\mathrm{ch,top}}\mathrm{\ compatibility})$
\textup{].}
For a holomorphic--topological theory $T$ with boundary chart algebra
$\cA$, write
\[
\mathcal O^{\mathrm{phys}}_{\mathrm{bulk}}(T)
\]
for the BRST complex of physical local bulk observables in the chosen
Costello--Gwilliam analytic completion.  The open--closed comparison
map is a chain map
\begin{equation}
\label{eq:master-oca-map-gate}
\mathrm{OCA}_T\colon
\mathcal O^{\mathrm{phys}}_{\mathrm{bulk}}(T)
\longrightarrow
Z^{\mathrm{der}}_{\mathrm{ch}}(\cA)
=C^\bullet_{\mathrm{ch}}(\cA,\cA)
\end{equation}
compatible with the closed-colour action on the boundary algebra.
If
\begin{equation}
\label{eq:master-oca-quasi-isomorphism-gate}
\mathrm{OCA}_T\ \text{is a quasi-isomorphism}
\end{equation}
then~\eqref{eq:master-oca-map-gate} identifies the BRST complex of
physical bulk observables with the derived chiral center of the
boundary chart.  In general, the two displayed complexes are joined by
the open--closed map~\eqref{eq:master-oca-map-gate}.
\end{proposition}

\begin{proposition}[HKR on an affine chart;
$\ClaimStatusProvedElsewhere$, physical realization
$\ClaimStatusConditional$]
\label{prop:master-oca-bmodel}
Let $A=k[x_1,\ldots,x_n]$.  The HKR antisymmetrization map
\[
 \operatorname{HKR}\colon
 \Lambda_A^p\operatorname{Der}_k(A)longrightarrow C^p(A,A)
\]
is given by
\[
 \operatorname{HKR}(X_1\wedge\cdots\wedge X_p)
 (a_1,\ldots,a_p)
 =\frac1{p!}\sum_{\sigma\in S_p}\!\operatorname{sgn}(\sigma)
 \prod_{j=1}^pX_{\sigma(j)}(a_j).
\]
It induces
\[
 A\otimes\Lambda^\bullet
 \langle\partial_{x_1},\ldots,\partial_{x_n}\rangle
 \xrightarrow{\sim}\mathrm{HH}^\bullet(A,A).
\]
For a physical theory $T$ with boundary chart $A$, a quasi-isomorphism
$\beta_T\colon\mathcal O^{\mathrm{phys}}_{\mathrm{bulk}}(T)	o
\Gamma(\mathbb A^n,\Lambda^\bullet T_{\mathbb A^n})$ makes
$\operatorname{HKR}\circ\beta_T$ the open--closed comparison.
\end{proposition}

\begin{proof}
The Hochschild--Kostant--Rosenberg theorem~\cite{HKR62}
gives, for every smooth commutative $k$-algebra $A$,
\[
\mathrm{HH}^\bullet(A, A)
 \;\cong\; \Lambda^\bullet_A \, \mathrm{Der}_k(A).
\]
For the polynomial ring,
$\operatorname{Der}_k(A)=\bigoplus_iA\partial_{x_i}$.  The displayed
antisymmetrization realizes the theorem.  Composition with
$\beta_T$ proves the conditional physical statement.
\end{proof}

\begin{remark}[Scope and extension]
\label{rem:master-oca-bmodel-scope}
The algebraic HKR identification extends to $X$ smooth affine over $k$ by the
same argument, with the trivialisation $T_X$ replaced by the
$\mathcal O_X$-module of derivations. For $X$ smooth projective, the
sheafy HKR \textup{(}polyvector fields = $\bigoplus_p
H^q(X, \Lambda^p T_X)$ in bigraded $(p, q)$\textup{)} computes the
Hochschild groups of $\mathrm{Perf}(X)$.  Compatibility of products
and brackets uses the Todd-twisted derived HKR map.  A physical
application additionally constructs
$\beta_T\colon\mathcal O^{\mathrm{phys}}_{\mathrm{bulk}}(T)\to
\mathrm R\Gamma(X,\Lambda^\bullet T_X)$ in the chosen BV
factorization model.
\end{remark}

\begin{lemma}[HKR for smooth affine and smooth projective schemes;
$\ClaimStatusProvedElsewhere$]
\label{lem:master-oca-sheafy-hkr-extension}
The OCA of Proposition~\ref{prop:master-oca-bmodel} extends:
\begin{enumerate}[label=\textup{(\arabic*)}]
\item \emph{Smooth affine $X$ over $k$.}
For $A = \mathcal O(X)$ the regular function ring of a smooth affine
variety $X$,
\[
\operatorname{HKR}_X\colon
\Gamma(X, \Lambda^\bullet T_X)
\;\xrightarrow{\;\sim\;}\;
\mathrm{HH}^\bullet(A, A),
\]
the HKR isomorphism for smooth commutative algebras.
\item \emph{Smooth projective $X$ over $k$.}
The category $\mathrm{Perf}(X)$ has sheafified HKR decomposition
\[
\operatorname{HKR}_X\colon
\mathrm R\Gamma\bigl(X, \Lambda^\bullet T_X\bigr)
\;\xrightarrow{\;\sim\;}\;
\mathrm{HH}^\bullet\bigl(\mathrm{Perf}(X)\bigr)
\;=\;
\bigoplus_{p+q = \bullet} H^q\bigl(X, \Lambda^p T_X\bigr).
\]
\end{enumerate}
\end{lemma}

\begin{proof}
\textup{(1)} For $X$ smooth affine, $A = \mathcal O(X)$ is a smooth
commutative $k$-algebra; the HKR theorem~\cite{HKR62}
\textup{(}\cite[\S 9.4.4]{Wei94}\textup{)} gives
$\mathrm{HH}^\bullet(A, A) \cong \Lambda^\bullet_A \mathrm{Der}_k(A)
\cong \Gamma(X, \Lambda^\bullet T_X)$. The argument of
Proposition~\ref{prop:master-oca-bmodel} for $X = \mathbb A^n$
applies verbatim with the polynomial-ring trivialisation
replaced by the locally free $\mathcal O_X$-module structure on
$T_X$.

\textup{(2)} For $X$ smooth projective, $D^b(\mathrm{Coh}(X))$ is
proper smooth and admits a Serre functor; its Hochschild cohomology
is computed by the sheafy HKR
\textup{(}see, e.g., the classical statement in
\cite[\S 9.4]{Wei94} extended to schemes via Swan-type sheaf
cohomology; modern treatment in
Calaque--Willwacher~\cite{CalaqueWillwacher21} for the higher-genus
operadic refinement\textup{).} The Hodge decomposition of polyvector
field sheaf cohomology gives the bigraded RHS.

This is the algebraic HKR statement.  A physical open--closed theorem
is obtained after constructing the bulk-to-polyvector map $\beta_T$
and its compatibility with the boundary action.
\end{proof}

\begin{corollary}[Hochschild cohomology of $K3$;
$\ClaimStatusProvedHere$]
\label{cor:master-oca-rozansky-witten-k3}
For a smooth projective K3 surface,
\[
\mathrm{HH}^p\bigl(\mathrm{Perf}(K3)\bigr) =
\begin{cases}
1 & p = 0,\,4,\\
22 & p = 2,\\
0 & \text{otherwise,}
\end{cases}
\qquad\sum_p \dim \mathrm{HH}^p = 24.
\]
If a Rozansky--Witten boundary condition supplies a quasi-isomorphism
$\beta_{\mathrm{RW}}$ from physical bulk observables to derived
polyvectors, then HKR identifies its physical bulk cohomology with
these groups.
\end{corollary}

\begin{proof}
By the sheafy HKR
\textup{(}Lemma~\ref{lem:master-oca-sheafy-hkr-extension}~\textup{(2)}\textup{)}
applied to $X = K3$:
\[
\mathrm{HH}^p(K3) \;=\; \bigoplus_{q + r = p} H^q(K3, \Lambda^r T_{K3}).
\]
$K3$ has trivial canonical bundle, so $\Lambda^2 T_{K3} \cong
K_{K3}^{\vee} \cong \mathcal O_{K3}$. The holomorphic symplectic form
$\sigma \in H^0(K3, \Omega^2)$ gives $T_{K3} \cong \Omega^1_{K3}$,
hence $H^q(K3, T_{K3}) = h^{1, q}(K3)$. The K3 Hodge diamond
$h^{0,0} = h^{2,0} = h^{0,2} = h^{2,2} = 1$, $h^{1,1} = 20$, with
all $h^{p, q}$ for $p + q$ odd equal to zero, yields:
\begin{align*}
\mathrm{HH}^0 &= H^0(K3, \mathcal O_{K3}) = 1,\\
\mathrm{HH}^1 &= H^0(T_{K3}) \oplus H^1(\mathcal O_{K3}) = 0 + 0 = 0,\\
\mathrm{HH}^2 &= H^0(\Lambda^2 T_{K3}) \oplus H^1(T_{K3}) \oplus
                H^2(\mathcal O_{K3}) = 1 + 20 + 1 = 22,\\
\mathrm{HH}^3 &= H^1(\Lambda^2 T_{K3}) \oplus H^2(T_{K3}) = 0 + 0 = 0,\\
\mathrm{HH}^4 &= H^2(\Lambda^2 T_{K3}) = 1.
\end{align*}
Total: $1 + 0 + 22 + 0 + 1 = 24 = \chi_{\mathrm{top}}(K3)$, agreeing
with the rank of the K3 Mukai cohomology lattice
$\Lambda_{\mathrm{Muk}}(K3) = H^*(K3, \mathbb Z)$ as expected from the
Mukai--C{\u a}ld{\u a}raru identity
$\dim \mathrm{HH}^*(\mathrm{Coh}(X)) = \dim H^*(X, \mathbb C)$ for
smooth projective $X$.

The physical interpretation is carried by the map
$\beta_{\mathrm{RW}}$ stated in the corollary.
\end{proof}

\begin{corollary}[Hochschild dimensions of K3 Hilbert schemes;
$\ClaimStatusProvedHere$]
\label{cor:master-oca-k3-hilbert-hierarchy}
For each $n \geq 1$, the Hilbert scheme $K3^{[n]}$ is a smooth
projective hyperk{\"a}hler $2n$-fold, and
$\dim \mathrm{HH}^*(\mathrm{Perf}(K3^{[n]})) = \dim H^*(K3^{[n]},
\mathbb C)$, computed by Göttsche~\cite{Gottsche1990}:
\[
\sum_{n \geq 0} q^n\, \dim H^*(K3^{[n]}, \mathbb C)
\;=\;
\prod_{m \geq 1}(1-q^m)^{-24}.
\]
The first several values:
$\dim H^*(K3^{[1]}) = 24$,
$\dim H^*(K3^{[2]}) = 324$,
$\dim H^*(K3^{[3]}) = 3200$.
A Rozansky--Witten bulk comparison for $K3^{[n]}$ transports these
groups to physical bulk cohomology.
\end{corollary}

\begin{proof}
Smoothness and hyperk{\"a}hler structure of $K3^{[n]}$ are standard
\textup{(}Beauville 1983\textup{).} The sheafy HKR
\textup{(}Lemma~\ref{lem:master-oca-sheafy-hkr-extension}~\textup{(2)}\textup{)}
gives $\dim \mathrm{HH}^*(K3^{[n]}) = \dim H^*(K3^{[n]}, \mathbb C)$
by Mukai--C{\u a}ld{\u a}raru. Göttsche's generating-function formula
gives the explicit dimensions \textup{(}derived from the symmetric
product / Hilbert-scheme decomposition and the K3 Hodge structure;
see \cite{Gottsche1990}\textup{).} Under the Rozansky--Witten bulk
comparison, these derived polyvector fields give the physical local
observables for the corresponding hyperk{\"a}hler target.
\end{proof}

\begin{remark}[Cross-link to K3$\times$E gravity-line residual]
\label{rem:master-oca-rw-k3-cross-link}
Corollary~\ref{cor:master-oca-rozansky-witten-k3} computes the
algebraic HKR target for the K3 factor of $K3 \times E$. The proposed
gravity-line algebra
of $K3 \times E$ involves additionally:
\begin{enumerate}[label=\textup{(\arabic*)}]
\item the elliptic fibre $E$ contribution
\textup{(}current envelope $\mathrm{Cur}_E$ in the
$(\beta, \delta, \varepsilon, \tau)$ datum of
Convention~\ref{conv:master-hall-borcherds-datum}\textup{);}
\item the Hall product encoding K3 wall-crossing
\textup{(}beyond the rigid HKR / Hodge identification used
here\textup{).}
\end{enumerate}
The HKR calculation above gives the finite-dimensional Hochschild
Gerstenhaber structure of the K3 factor.  The Hall correspondences
required for the K3-BKM comparison are carried by the maps
$(\delta,\varepsilon)$ of
Remark~\ref{rem:master-witness-K3E-ramification}.
\end{remark}

\begin{corollary}[Hochschild cohomology of an elliptic curve;
$\ClaimStatusProvedHere$]
\label{cor:master-oca-bmodel-elliptic}
For a smooth elliptic curve $E$,
\[
\mathrm{HH}^p\bigl(D^b(\mathrm{Coh}(E))\bigr) =
\begin{cases}
1 & p = 0,\,2,\\
2 & p = 1,\\
0 & \text{otherwise,}
\end{cases}
\qquad \sum_p \dim \mathrm{HH}^p = 4.
\]
A bulk-to-polyvector quasi-isomorphism transports this calculation to
the corresponding physical theory.
\end{corollary}

\begin{proof}
$E$ has trivial canonical bundle and $T_E \cong \mathcal O_E$
\textup{(}translation-invariant frame on an abelian variety\textup{).}
Since $E$ is a complex $1$-fold, $\Lambda^p T_E = 0$ for $p \geq 2$.
The Hodge diamond
$h^{0,0} = h^{1,0} = h^{0,1} = h^{1,1} = 1$ \textup{(}all others
zero\textup{)} gives
$H^q(E, T_E) = H^q(E, \mathcal O_E) =
\begin{cases} 1 & q = 0, 1\\ 0 & q \geq 2\end{cases}$.
The sheafy HKR
\textup{(}Lemma~\ref{lem:master-oca-sheafy-hkr-extension}~\textup{(2)}\textup{)}
yields:
\begin{align*}
\mathrm{HH}^0 &= H^0(\mathcal O_E) = 1,\\
\mathrm{HH}^1 &= H^0(T_E) \oplus H^1(\mathcal O_E) = 1 + 1 = 2,\\
\mathrm{HH}^2 &= H^1(T_E) = 1,
\end{align*}
with $\mathrm{HH}^p = 0$ for $p \geq 3$. Total
$\sum_p \dim \mathrm{HH}^p = 4 = b_0(E) + b_1(E) + b_2(E) =
\dim H^*(E, \mathbb C)$, in accord with the Mukai--C{\u a}ld{\u a}raru
identity.
\end{proof}

\begin{remark}[Künneth: abelian Hochschild of $K3 \times E$]
\label{rem:master-oca-k3-e-kunneth}
By the K{\"u}nneth formula for derived Hochschild cohomology of
smooth proper products,
\[
\mathrm{HH}^\bullet\bigl(D^b(\mathrm{Coh}(K3 \times E))\bigr)
\;=\;
\mathrm{HH}^\bullet\bigl(D^b(\mathrm{Coh}(K3))\bigr) \otimes
\mathrm{HH}^\bullet\bigl(D^b(\mathrm{Coh}(E))\bigr),
\]
with total dimension $24 \cdot 4 = 96 = \dim H^*(K3, \mathbb C)
\cdot \dim H^*(E, \mathbb C) = \dim H^*(K3 \times E, \mathbb C)$.
This $96$-dimensional structure is the \emph{abelian Hochschild}
content of $K3 \times E$ \textup{(}commutative associative Hochschild
algebra structure via cup product on polyvector fields, with Schouten
bracket as the closed-colour
$\mathsf{SC}^{\mathrm{ch,top}}$-brace per
Corollary~\ref{cor:master-oca-bmodel-schouten}\textup{).}

The K3-BKM algebra $\mathfrak g_{\Delta_5}$ that appears in the
$K3 \times E$ gravity-line setup is \emph{infinite-dimensional}
and graded by the Lorentzian root lattice of the $\Delta_5$
denominator, with root multiplicities given by the coefficients of
$\phi_{0,1}$,
whereas the Hochschild space above has dimension~$96$ and carries the
polyvector cup product and Schouten bracket.  The Hall enhancement
$Y^+_{\mathrm{Hall}}(K3)$ is formed by correspondences on the moduli
stack of K3 stable sheaves.  The maps $(\delta,\varepsilon)$ of
Convention~\ref{conv:master-hall-borcherds-datum} compare its
root-graded Hall operations with the K3-BKM gravity-line algebra.
\end{remark}

\begin{corollary}[Hochschild cohomology of the quintic threefold;
$\ClaimStatusProvedHere$]
\label{cor:master-oca-quintic}
For $Y = Y_5$ the quintic Calabi--Yau threefold in $\mathbb P^4$, the
Hochschild cohomology of $\mathrm{Perf}(Y_5)$ has dimension vector
\[
\bigl(\dim \mathrm{HH}^p(D^b(\mathrm{Coh}(Y_5)))\bigr)_{p = 0,\ldots,6}
\;=\;
(1,\, 0,\, 101,\, 4,\, 101,\, 0,\, 1),
\qquad
\sum_p \dim \mathrm{HH}^p = 208.
\]
\end{corollary}

\begin{proof}
The quintic is a smooth Calabi--Yau $3$-fold with Hodge numbers
$h^{0, 0} = h^{3, 3} = 1$, $h^{1, 1} = 1$, $h^{2, 1} = h^{1, 2} = 101$,
$h^{3, 0} = h^{0, 3} = 1$, $h^{2, 2} = 1$, other Hodge numbers zero.
For a CY $3$-fold, the holomorphic volume form provides
$\Lambda^p T_Y \cong \Omega^{3-p}_Y$ \textup{(}via
$K_Y^\vee = \det T_Y \cong \mathcal O_Y$\textup{).} Hence
$H^q(\Lambda^p T_Y) = h^{3 - p,\, q}$.

By sheafy HKR
\textup{(}Lemma~\ref{lem:master-oca-sheafy-hkr-extension}~\textup{(2)}\textup{)}:
\begin{align*}
\mathrm{HH}^0 &= h^{3, 0} = 1,\\
\mathrm{HH}^1 &= h^{2, 0} + h^{3, 1} = 0 + 0 = 0,\\
\mathrm{HH}^2 &= h^{1, 0} + h^{2, 1} + h^{3, 2} = 0 + 101 + 0 = 101,\\
\mathrm{HH}^3 &= h^{0, 0} + h^{1, 1} + h^{2, 2} + h^{3, 3}
               = 1 + 1 + 1 + 1 = 4,\\
\mathrm{HH}^4 &= h^{0, 1} + h^{1, 2} + h^{2, 3} = 0 + 101 + 0 = 101,\\
\mathrm{HH}^5 &= h^{0, 2} + h^{1, 3} = 0 + 0 = 0,\\
\mathrm{HH}^6 &= h^{0, 3} = 1.
\end{align*}
Total $1 + 0 + 101 + 4 + 101 + 0 + 1 = 208 = \dim H^*(Y_5, \mathbb C)$,
consistent with the Mukai--C{\u a}ld{\u a}raru identity.

The four one-dimensional summands of $\mathrm{HH}^3$ are
$H^0(\mathcal O_Y)$, $H^1(\Omega_Y^1)$,
$H^2(\Omega_Y^2)$, and $H^3(\Omega_Y^3)$.  The
$101$-dimensional terms in degrees $2$ and $4$ arise from
$H^1(\Omega_Y^2)$ and $H^2(\Omega_Y^1)$, respectively.
\end{proof}

\begin{remark}[Mirror exchange visible in OCA]
\label{rem:master-oca-mirror-quintic}
Mirror symmetry relates the quintic $Y_5$ and the mirror quintic
$\tilde Y_5$, exchanging $h^{1, 1} \leftrightarrow h^{2, 1}$:
$\tilde h^{1, 1} = 101$, $\tilde h^{2, 1} = 1$. At the level of
Hochschild cohomology,
$\bigl(\dim \mathrm{HH}^p(\tilde Y_5)\bigr) =
(1,0,1,204,1,0,1).
Indeed, the middle term is
\[
 \dim\mathrm{HH}^3(\tilde Y_5)
 =h^{0,0}+h^{1,1}+h^{2,2}+h^{3,3}
 =1+101+101+1=204.
\]
Thus $\dim\mathrm{HH}^*(\tilde Y_5)=208$, equal to the total Betti
number of the mirror.  The exchange
$h^{1,1}\leftrightarrow h^{2,1}$ moves the two $101$-dimensional
summands from degrees $2$ and $4$ for $Y_5$ to degree~$3$ for
$\tilde Y_5$.  This is the HKR form of the exchange between Kähler
and complex-structure directions.
\end{remark}

\begin{corollary}[Schouten--Nijenhuis bracket as closed-colour
$\mathsf{SC}^{\mathrm{ch,top}}$-brace;
$\ClaimStatusProvedElsewhere$, physical realization
$\ClaimStatusConditional$]
\label{cor:master-oca-bmodel-schouten}
Under the OCA identification of
Proposition~\ref{prop:master-oca-bmodel}, the closed-colour
$\mathsf{SC}^{\mathrm{ch,top}}$-brace bracket on
$Z^{\mathrm{der}}_{\mathrm{ch}}(A) = \mathrm{HH}^\bullet(A, A)$
corresponds, for $A = k[x_1, \ldots, x_n]$, to the
Schouten--Nijenhuis bracket on polyvector fields
$\Lambda^\bullet T_{\mathbb A^n}$:
\[
[-, -]_{\mathrm{brace}}\bigl|_{\mathrm{HH}^\bullet}
 \;\xleftrightarrow{\;\mathrm{HKR}\;}\;
[-, -]_{SN}\bigl|_{\Lambda^\bullet T_{\mathbb A^n}}.
\]
Consequently a physical open--closed comparison compatible with
braces transports this Gerstenhaber structure to the bulk operator
complex.
\end{corollary}

\begin{proof}
The Gerstenhaber bracket on $\mathrm{HH}^\bullet(A, A)$ for any
associative algebra $A$ was introduced
in~\cite{Ger63}; it gives a degree-$-1$ graded Lie bracket compatible
with the cup product, exhibiting $\mathrm{HH}^\bullet(A, A)$ as a
Gerstenhaber algebra. Kontsevich~\cite{Kon99} promotes this to a
brace algebra and ultimately to an $E_2$-algebra structure compatible
with the $\mathsf{SC}^{\mathrm{ch,top}}$ closed colour.

For $A=k[x_1,\ldots,x_n]$, the HKR map of
Proposition~\ref{prop:master-oca-bmodel} identifies Hochschild classes
with polyvector fields.  The Schouten bracket is determined by
\[
 [f,g]_{SN}=0,
 \qquad
 [X,f]_{SN}=X(f),
 \qquad
 [X,Y]_{SN}=[X,Y]_{\mathrm{Lie}},
\]
and the graded biderivation rule
\[
 [P,Q\wedge R]_{SN}
 =[P,Q]_{SN}\wedge R
 +(-1)^{(|P|-1)|Q|}Q\wedge[P,R]_{SN}.
\]
HKR intertwines this bracket with the Gerstenhaber bracket on
cohomology; see~\cite[\S9.4]{Wei94}.
\end{proof}

\begin{lemma}[A fibre of the graded Euler series;
$\ClaimStatusProvedHere$]
\label{lem:master-scalar-non-faithfulness}
The graded Euler series
\[
\chi\colon \operatorname{Ob}(\mathrm{FiltCh}_k)
\longrightarrow \mathbb Z(\!(q)\!)
\]
has fibres containing several quasi-isomorphism classes.  Consequently
a Hall product, Drinfeld pairing, current-envelope locality structure,
derived-centre trace, or
$\mathsf{SC}^{\mathrm{ch,top}}$-action requires its own morphisms and
identities in addition to the value of $\chi$.
\end{lemma}

\begin{proof}
Place every vector space in filtration weight~$0$.  Let
\[
 C=k[0],\qquad
 K=\bigl(k[0]\oplus k[1],d=0\bigr),\qquad
 C'=C\oplus K.
\]
Then $\chi(K)=1-1=0$, hence $\chi(C')=\chi(C)=1$.  Their homology
dimensions are
\[
 \dim H^\bullet(C)=1,
 \qquad
 \dim H^\bullet(C')=3.
\]
Thus $C$ and $C'$ represent distinct quasi-isomorphism classes in one
Euler-series fibre.  Additional algebraic operations contain still
finer data.
\end{proof}

\begin{theorem}[A rank-$16$ graded-character fibre;
$\ClaimStatusProvedHere$]
\label{thm:master-scalar-nonfaithful-witness-c16}
The lattice vertex operator algebras
\[
\cA_1 \;=\; \Vlat_{E_8 \oplus E_8},
\qquad
\cA_2 \;=\; \Vlat_{D_{16}^+}
\;=\; \Vlat_{\mathrm{Spin}(32)/\mathbb Z_2},
\]
attached to the two even unimodular positive-definite lattices of
rank~$16$ \textup{(}\cite{Bor86, FLM88}\textup{)} satisfy
\begin{enumerate}[label=\textup{(\roman*)}]
\item the same central charge $c(\cA_1) = c(\cA_2) = 16$;
\item the same graded character
\[
\chi(\cA_1)(q) \;=\; \chi(\cA_2)(q)
\;=\; \frac{E_8(q)}{\eta(q)^{16}},
\qquad
E_8(q) = 1 + 480\sum_{n\geq 1}\sigma_7(n)\,q^n;
\]
\item distinct weight-one Lie algebras
$\mathfrak e_8\oplus\mathfrak e_8$ and $\mathfrak d_{16}$, which
distinguish their chiral algebra structures and their
$\mathsf{SC}^{\mathrm{ch,top}}$-objects on the open colour.
\end{enumerate}
This is the concrete realisation of
Lemma~\ref{lem:master-scalar-non-faithfulness} at the smallest
central charge where the standard landscape exhibits two distinct
chiral algebras with equal graded character.
\end{theorem}

\begin{proof}
\textup{(}i\textup{)} For a lattice VOA $\Vlat_L$ with $L$ even
integral of rank $r$, the central charge is $c(\Vlat_L) = r$. Both
lattices have rank $16$.

\textup{(}ii\textup{)} The graded character is
$\chi(\Vlat_L)(q) = \theta_L(q) / \eta(q)^r$. For $L$ even unimodular
positive-definite of rank $r$ divisible by $8$, the theta function
$\theta_L$ is a holomorphic modular form of weight $r/2$ on
$\mathrm{SL}_2(\mathbb Z)$. At $r = 16$ the space
$M_8(\mathrm{SL}_2(\mathbb Z))$ is one-dimensional, spanned by the
Eisenstein series $E_8(q)$ \textup{(}cusp forms start at weight
$12$\textup{)}; both $\theta_{E_8 \oplus E_8}$ and
$\theta_{D_{16}^+}$ have constant term~$1$, hence
$\theta_{E_8 \oplus E_8} = \theta_{D_{16}^+} = E_8(q)$.

\textup{(}iii\textup{)} The weight-$1$ subspaces are the affine
Lie algebras attached to the norm-$2$ vectors of $L$:
\[
(\cA_1)_1 \;=\; \widehat{\mathfrak e_8 \oplus \mathfrak e_8}_1,
\qquad
(\cA_2)_1 \;=\; \widehat{\mathfrak{so}(32)}_1
                                 = \widehat{\mathfrak d_{16}}_1,
\]
both of dimension $496$, both of rank $16$, both with level-$1$
Sugawara central charge $16$ \textup{(}consistent with
\textup{(}ii\textup{)}\textup{)}. The horizontal Lie algebras have
different decomposition types: $\mathfrak e_8 \oplus \mathfrak e_8$
has two simple summands of rank $8$ each, whereas $\mathfrak d_{16}$
is simple of rank $16$. Equivalently, their Dynkin diagrams differ
\textup{(}two disconnected copies of $E_8$ versus the connected
$D_{16}$ fork\textup{),} and so do the multisets of fundamental
degrees: $\{2,8,12,14,18,20,24,30\}^{\sqcup 2}$ for
$\mathfrak e_8 \oplus \mathfrak e_8$ versus
$\{2,4,6,\ldots,28,30,16\}$ for $\mathfrak d_{16}$ \textup{(}where
the extra $16$ is the Pfaffian degree\textup{).} The OPEs on
$(\cA)_1\otimes(\cA)_1$ therefore recover distinct affine Lie
algebras.  This distinguishes $\cA_1$ from $\cA_2$ as chiral algebras
and as $\mathsf{SC}^{\mathrm{ch,top}}$-objects on the open colour.
\end{proof}

\begin{remark}[K3$\times$E ramification: the load-bearing application;
construction~\textup{(2)} $\ClaimStatusConditional$]
\label{rem:master-witness-K3E-ramification}
The $c=16$ example of
Theorem~\ref{thm:master-scalar-nonfaithful-witness-c16} exhibits two
chiral algebras in one graded-character fibre.  The K3$\times E$ case
uses this fibre phenomenon.

The scalar identity $Z_{\mathrm{BPS}}^{K3 \times E} =
(\Phi_{10}^{\mathrm{un}})^{-1} = \Delta_5^{-2}$ specifies a graded
character at level $4$. Two natural chain-level candidate carriers
share this character:
\begin{enumerate}[label=\textup{(\arabic*)}]
\item the canonical Borcherds envelope
$\mathbf H_{\Delta_5}$, constructed from the BKM algebra
$\mathfrak g_{\Delta_5}$ via the additive lift $\Delta_5
= \mathrm{Grit}(\eta^9 \vartheta_1) \in S_5(K(1))$
\textup{(}\cite[Thm.~6.1]{Gritsenko1999}\textup{);}
\item the Schiffmann--Vasserot Hall envelope
$\mathrm{Cur}_E\bigl(D(Y^+_{\mathrm{Hall}}(K3))\bigr)$
\textup{(}\cite{SMMV23}; chain-level construction conditional on the
completion theorem of
Convention~\ref{conv:master-hall-borcherds-datum}\textup{).}
\end{enumerate}
The four-part datum $(\beta,\delta,\varepsilon,\tau)$ defines the
comparison from~\textup{(2)} to~\textup{(1)}.  At the present status,
the two chain complexes lie in the same scalar-character fibre.
Lemma~\ref{lem:master-scalar-non-faithfulness} shows that this fibre
contains several chiral-algebra equivalence classes.  Thus an
equivalence
\[
 \mathbf H_{\Delta_5}\simeq
 \mathrm{Cur}_E\bigl(D(Y^+_{\mathrm{Hall}}(K3))\bigr)
\]
is the precise theorem governed by the four comparison maps.

The map $\beta$ specifies the positive-half action on
$U^{\mathrm{ch}}(\mathfrak n_+(\mathfrak g_{\Delta_5}))$.  The maps
$\delta$ and $\varepsilon$ add the Drinfeld pairing of the Hall
algebra and current locality along $E$.  Together these are the data
required to compare K3 Donaldson--Thomas wall-crossing
\textup{(}\cite{KontsevichSoibelman}\textup{)} with the proposed
gravity-line object.
\end{remark}

\begin{remark}
\label{rem:master-scalar-shadow-promotion}
The scalar character identity at level~$5$
\textup{(}e.g.\ the denominator
$Z_{\mathrm{BPS}}^{K3\times E}(Z)
 = (\Phi_{10}^{\mathrm{un}}(Z))^{-1}
 = \Delta_5(Z)^{-2}$\textup{)} lies in the scalar target.  The
four-part datum $(\beta,\delta,\varepsilon,\tau)$ of
Convention~\ref{conv:master-hall-borcherds-datum} supplies the
comparison with a level-$4$ gravity-line algebra.
\end{remark}

\begin{convention}[Hall--Borcherds gravity-line chain datum
at level $5{\to}4$; \ClaimStatusDefinitional]
\label{conv:master-hall-borcherds-datum}
A chain-level gravity-line comparison at level $5 \to 4$ for the
target $X = K3 \times E$ is the four-part datum
$(\beta, \delta, \varepsilon, \tau)$:
\begin{enumerate}[label=\textup{(\arabic*)}]
\item $\beta\colon Y^+_{\mathrm{Hall}}(K3 \times E)
\to U^{\mathrm{ch}}(\mathfrak n_+(\mathfrak g_{\Delta_5}))$,
an $E_1$-chiral bialgebra morphism on the reference elliptic curve
\textup{(}positive-half Hall--Borcherds comparison\textup{)};
\item $\delta\colon D(Y^+_{\mathrm{Hall}})
\to D(U^{\mathrm{ch}}(\mathfrak n_+))$,
the completed Drinfeld-double extension, compatible with the Hall
pairing, the PBW filtration, and the integral form;
\item $\varepsilon\colon \mathrm{Cur}_E\, D(Y^+_{\mathrm{Hall}})
\to \mathrm{Op}^{\mathrm{grav}}_{K3 \times E}$,
the current-envelope morphism along the elliptic fibre, landing in the
gravitational boundary line algebra
\textup{(}filtered $\mathsf{SC}^{\mathrm{ch,top}}$-morphism\textup{)};
\item $\tau$, the derived-centre trace identity
\[
\Tr_{Z^{\mathrm{der}}_{\mathrm{ch}}}(\varepsilon)
 = (\Phi_{10}^{\mathrm{un}})^{-1} = \Delta_5^{-2}.
\]
\end{enumerate}
The comparison is formed through finite truncations
$W=(\text{height},\text{charge},\text{weight})$.  For each $W$ one
constructs maps $\beta_W$, $\delta_W$, $\varepsilon_W$, and
$\tau_W$ and proves the following eight assertions:
\begin{enumerate}[label=\textup{(C\arabic*)}]
\item the finite Hall source exists;
\item the root multiplicities match $\Delta_5$;
\item the associated graded of the PBW filtration has the prescribed
root spaces;
\item parity signs match super-root multiplicities;
\item the Hall pairing is non-degenerate after completion;
\item the Drinfeld double respects current locality on $E$;
\item the $\mathsf{SC}^{\mathrm{ch,top}}$ mixed operations commute
with the trace;
\item the inverse limit over $W$ is Mittag--Leffler, or otherwise
controlled.
\end{enumerate}
These assertions produce an inverse system of chain-level comparison
maps.  Its limit gives $(\beta,\delta,\varepsilon,\tau)$ under the
Mittag--Leffler clause.

Gritsenko--Nikulin construct the automorphic product
\cite[Example~2.4, equation~(2.16)]{GritsenkoNikulin1998}
\[
 \Delta_5(Z)
 =(qrs)^{1/2}
 \prod_{(n,l,m)>0}
 \bigl(1-q^nr^ls^m\bigr)^{f_1(nm,l)},
\]
where $f_1(n,l)$ are the Fourier coefficients of the weak Jacobi form
$\phi_{0,1}$.  Their Section~5.1.1 interprets this product as an
automorphic correction of a rank-three Lorentzian Kac--Moody system.
This determines the denominator and the root multiplicities.  A Lie
bialgebra cobracket, its quantization, the Hall comparison, and the
current-envelope map are the additional structures represented by
$\beta$, $\delta$, and $\varepsilon$ above; their present status is
$\ClaimStatusConjectured$.
\end{convention}

\begin{lemma}[Classical double of the zero abelian Lie bialgebra;
$\ClaimStatusProvedHere$]
\label{lem:master-delta-cartan-window}
Let $\mathfrak h$ be a finite-dimensional vector space with the
abelian Lie bracket and zero cobracket.  Its classical Drinfeld double
is
\[
D(\mathfrak h)=\mathfrak h\oplus\mathfrak h^*.
\]
Its Lie bracket is zero, and its invariant bilinear form is
\[
 \langle x+\xi,y+\eta\rangle_D
 =\xi(y)+\eta(x).
\]
Thus $(D(\mathfrak h),\mathfrak h,\mathfrak h^*)$ is the Manin triple
associated with the zero abelian Lie bialgebra.  The evaluation pairing
is a bilinear form on the double.
\end{lemma}

\begin{proof}
The bracket on $\mathfrak h^*$ is dual to the cobracket on
$\mathfrak h$, and the mixed bracket is the sum of the two coadjoint
actions.  The zero bracket and zero cobracket make all three terms
zero.  The evaluation form is non-degenerate, invariant, and restricts
to zero on each summand; hence the two summands are Lagrangian.  This
is the standard double construction
\cite[\S4]{EtingofSchiffmann98}.
\end{proof}

\begin{lemma}[Heisenberg current algebra from a bilinear form;
$\ClaimStatusProvedHere$]
\label{lem:master-epsilon-cartan-window}
Let $\mathfrak h$ be a finite-dimensional vector space and
$\kappa\colon\mathfrak h\otimes\mathfrak h\to k$ a symmetric
bilinear form.  On a formal disc, the central extension
\[
 \widehat{\mathfrak h((t))}_\kappa
 =\mathfrak h((t))\oplus kK
\]
has bracket
\[
 [x\otimes f,y\otimes g]
 =\kappa(x,y)\operatorname{Res}(f\,dg)K.
\]
Its chiral envelope on a smooth curve $E$ is the Heisenberg chiral
algebra $\mathcal H_{\mathfrak h,\kappa,E}$ with OPE
\[
 J^x(z)J^y(w)\sim\frac{\kappa(x,y)}{(z-w)^2}.
\]
\end{lemma}

\begin{proof}
Antisymmetry follows from
$\operatorname{Res}(f\,dg)=-\operatorname{Res}(g\,df)$, and the
centrality of $K$ gives the Jacobi identity.  Chiral induction from
this local Lie algebra gives the stated current algebra
\cite[\S2.5]{FrenkelBenZvi}.  An application to a Cartan subalgebra of
$\mathfrak g_{\Delta_5}$ therefore begins with an explicitly chosen
form $\kappa$ and a comparison from the Hall Cartan to
$\mathcal H_{\mathfrak h,\kappa,E}$.
\end{proof}

\begin{lemma}[Heisenberg character and the $\Delta_5^{-2}$ comparison;
$\ClaimStatusConditional$]
\label{lem:master-tau-cartan-window}
The vacuum character of a rank-$r$ Heisenberg chiral algebra is
\[
 \chi_{\mathcal H_r}(\tau)=\eta(\tau)^{-r}.
\]
For a positive-definite even lattice $L$ of rank $r$, its lattice
vertex algebra has character $\theta_L(\tau)/\eta(\tau)^r$.  The
Lorentzian rank-three root lattice of the $\Delta_5$ denominator lies
in a different analytic domain: its theta series requires a specified
regularization and convergence theorem.

The exact scalar target is
\[
 \Delta_5(Z)^{-2}
 =(qrs)^{-1}
 \prod_{(n,l,m)>0}
 \bigl(1-q^nr^ls^m\bigr)^{-2f_1(nm,l)}.
\]
The trace comparison asks for a graded chain object whose superdimension
in degree $(n,l,m)$ is $-2f_1(nm,l)$, together with the sewing map
$\tau$ that identifies its completed trace with this product.
\end{lemma}

\begin{proof}
The Heisenberg and positive-definite lattice character formulae are
standard \cite[\S2]{FrenkelBenZvi}.  Squaring the reciprocal of
Gritsenko--Nikulin's product
\cite[Example~2.4, equation~(2.16)]{GritsenkoNikulin1998}
gives the displayed three-variable product.  The construction of
$\tau$ is precisely the chain-level trace problem in
Convention~\ref{conv:master-hall-borcherds-datum}.
\end{proof}

\begin{remark}[Established and conditional components of
$(\beta,\delta,\varepsilon,\tau)$]
\label{rem:master-hb-cartan-window-synthesis}
The preceding results separate four mathematical statements:
\begin{itemize}[noitemsep]
\item Gritsenko--Nikulin's product determines the automorphic
denominator and the integers $f_1(nm,l)$.
\item A chosen zero abelian Lie bialgebra has the double of
Lemma~\ref{lem:master-delta-cartan-window}.
\item A chosen symmetric form on a Cartan space gives the Heisenberg
current algebra of Lemma~\ref{lem:master-epsilon-cartan-window}.
\item The equality of a completed chain-level trace with
$\Delta_5^{-2}$ is the conditional map $\tau$ of
Lemma~\ref{lem:master-tau-cartan-window}.
\end{itemize}
The Hall--Borcherds theorem therefore asks for a specified Lie
bialgebra on the Hall side, its completed Drinfeld double, a chiral
current realization on $E$, and a trace map recovering the exact
three-variable product.
\end{remark}

\begin{convention}[Chain data for a topological chiral enhancement;
\ClaimStatusDefinitional]
\label{conv:master-topologisation}
Let $V$ be a chiral algebra with conformal vector $T$.  A
\emph{topological chiral enhancement} is specified in a complete
filtered chain model by the following data:
\begin{enumerate}[label=\textup{(T\arabic*)}]
\item a degree-one derivation $Q$ with $Q^2=0$, an odd field $G$, and
the identity $[Q,G]=T$;
\item a complete exhaustive filtration preserved by $Q$, $G$, and the
chiral operations;
\item continuity of the operations for the chosen completion and a
finite-propagation estimate on every bounded-weight quotient;
\item an action of the stated topological operad together with a
homotopy-transfer comparison for two choices of contraction;
\item a null-homotopy of the anomaly cocycle in the deformation
complex of that operad action.
\end{enumerate}
The identity in T1 makes holomorphic translation $Q$-exact on
cohomology.  An $E_3^{\mathrm{top}}$ conclusion uses the operadic
action and comparison homotopies in T4.  Each theorem below states
exactly which entries of this package it constructs.
\end{convention}

\begin{theorem}[Bosonic-string BRST identities at matter central
charge $26$; $\ClaimStatusProvedElsewhere$]
\label{thm:master-vir-26-bc-brst-topologisation}
At matter central charge $c_m = 26$, the Virasoro algebra
$\mathrm{Vir}_{26}$ tensored with the conformal-weight-$(2, -1)$
fermionic ghost system $(b, c)$ of central charge $c_{bc} = -26$
carries the BRST current
\[
T_{bc} = -2{:}b\partial c{:} - {:}\partial b\, c{:},
\qquad
T_{\mathrm{tot}} = T_m + T_{bc},
\]
\[
j_{\mathrm{BRST}}
= {:}c(T_m+\tfrac12T_{bc}){:}+\tfrac32\partial^2c,
\qquad
Q_{\mathrm{BRST}}=\oint\frac{dz}{2\pi i}\,j_{\mathrm{BRST}}(z),
\qquad
G = b.
\]
Then
\[
Q_{\mathrm{BRST}}^2=0,
\qquad
[Q_{\mathrm{BRST}}, b(z)] = T_{\mathrm{tot}}(z).
\]
These identities give T1 and the BRST anomaly cancellation entering
T5 of Convention~\ref{conv:master-topologisation}.  A chosen complete
filtered model and a chain action of the desired topological operad
supply T2--T4.
\end{theorem}

\begin{proof}
The central charge of a fermionic $(b, c)$ system of weight $(h, 1-h)$
is $c_{bc}(h) = -2(6h^2 - 6h + 1)$
\textup{(}Polchinski~\cite[eq.~(2.5.12)]{Pol98}\textup{);} at $h = 2$
this gives $c_{bc} = -26$.

The standard BRST OPE calculation gives
$Q_{\mathrm{BRST}}^2=(c_m-26)\mathfrak a_{\mathrm{Vir}}$, where
$\mathfrak a_{\mathrm{Vir}}$ is the ghost-number-two Virasoro anomaly
operator; see~\cite[\S4.3]{Pol98}.  Its coefficient vanishes at
$c_m=26$.

The identity $[Q_{\mathrm{BRST}}, b(z)] = T_{\mathrm{tot}}(z)$ follows
from
\[
j_{\mathrm{BRST}}(z) b(w)
 \;\sim\; \frac{T_{\mathrm{tot}}(w)}{z - w} + \cdots,
\]
which follows from $b(z)c(w)\sim(z-w)^{-1}$ and the standard
stress-tensor OPEs.  The residue of the first-order pole gives the
displayed commutator.
\end{proof}

\begin{remark}[The Virasoro anomaly away from the critical point]
\label{rem:master-vir-noncritical-completed}
For the reparametrization BRST operator of
Theorem~\ref{thm:master-vir-26-bc-brst-topologisation}, the anomaly
class is
\[
 [Q_{\mathrm{BRST}}^2]
 =(c_m-26)[\mathfrak a_{\mathrm{Vir}}].
\]
The critical locus is therefore $c_m=26$.  Completion transports this
class to the completed deformation complex with the same coefficient.
Felder screening differentials form a separate family of complexes,
described in Corollary~\ref{cor:master-screening-brst-c-lt-minus-1}.
\end{remark}

\begin{theorem}[Conformal-weight constraint for the
reparametrization BRST current; $\ClaimStatusProvedHere$]
\label{thm:master-generalized-matter-ghost-brst}
Let $T_m$ have conformal weight~$2$, and let $bc(h,1-h)$ be a
fermionic ghost system with
\[
 \operatorname{wt}(b)=h,
 \qquad
 \operatorname{wt}(c)=1-h.
\]
Conformal covariance of a reparametrization BRST current with leading
term ${:}cT_m{:}$ forces
\[
 \operatorname{wt}({:}cT_m{:})=1,
 \qquad h=2,
 \qquad (\operatorname{wt}(b),\operatorname{wt}(c))=(2,-1).
\]
For this ghost system $c_{bc}=-26$, and the BRST anomaly coefficient
is $c_m-26$.  Thus the reparametrization construction of
Theorem~\ref{thm:master-vir-26-bc-brst-topologisation} is selected by
the simultaneous weight and anomaly equations.
\end{theorem}

\begin{proof}
The field ${:}cT_m{:}$ has weight $(1-h)+2=3-h$; setting this equal
to~$1$ gives $h=2$.  The central-charge formula
\[
 c_{bc}(h)=-2(6h^2-6h+1)
\]
then gives $c_{bc}(2)=-26$
\cite[equation~(2.5.12)]{Pol98}.  The standard BRST OPE gives anomaly
coefficient $c_m+c_{bc}=c_m-26$.
\end{proof}

\begin{remark}[Three distinct BRST constructions]
\label{rem:master-generalized-brst-coverage}
The reparametrization BRST complex uses the $(2,-1)$ ghost system and
the critical equation $c_m=26$.  Felder complexes use screening
operators between Coulomb-gas Fock modules.  An $N=2$ topological
twist uses one supercurrent as $Q$ and the other as $G$.  These three
constructions have distinct chain complexes and distinct nilpotency
identities.  A higher-spin $\mathcal W$-algebra uses the full BRST
master equation of
Corollary~\ref{cor:master-per-spin-antighost-linear-Winf}.
\end{remark}

\begin{corollary}[The higher-spin BRST master equation;
$\ClaimStatusConditional$]
\label{cor:master-per-spin-antighost-linear-Winf}
Let $\mathcal W$ be generated by fields $W^{(s)}$ of conformal weight
$s$ with specified OPE structure constants.  Introduce ghosts
$b_s,c_s$ of weights $(s,1-s)$.  A higher-spin BRST charge has the
form
\[
 Q_{\mathcal W}
 =\oint\left(
 \sum_s c_sW^{(s)}
 +\sum_{r,s,t}C_{rs}^{t}{:}c_rc_sb_t{:}
 +\text{higher ghost terms}
 \right).
\]
The coefficients and higher terms solve the master equation
$Q_{\mathcal W}^2=0$.  Under this equation,
\[
 [Q_{\mathcal W},b_s]=W^{(s)}_{\mathrm{tot}}
\]
is the per-spin antighost identity.  For
$\mathcal W_\infty[\lambda]$, the mixed OPEs among different spins
enter the quadratic and higher ghost terms.  A proof consists of an
explicit solution of this master equation in a stated completion.
\end{corollary}

\begin{proof}
The coefficient of each ghost monomial in $Q_{\mathcal W}^2$ is the
Chevalley--Eilenberg relation for the corresponding OPE bracket,
together with its quantum anomaly.  Solving these equations gives a
differential.  Contracting the BRST current with $b_s$ extracts the
constraint paired with $c_s$ and its ghost correction, which is
$W^{(s)}_{\mathrm{tot}}$.
\end{proof}

\begin{corollary}[Felder screening complexes;
$\ClaimStatusProvedElsewhere$]
\label{cor:master-screening-brst-c-lt-minus-1}
For coprime integers $p,q\geq2$, the Coulomb-gas realization at
\[
 c_{p,q}=1-6\frac{(p-q)^2}{pq}
\]
carries screening operators
\[
 S_\pm=\oint\frac{dz}{2\pi i}
 {:}e^{\alpha_\pm\phi(z)}{:}.
\]
Suitable powers of these operators give the differentials in Felder's
complexes of Fock modules, and the resulting cohomology realizes the
irreducible Virasoro minimal-model modules; see
\cite[\S15]{FrenkelBenZvi}.  This is a resolution theorem.  A
topological chiral enhancement adds a field $G$ and the identity
$[Q,G]=T$ of Convention~\ref{conv:master-topologisation} for the
chosen differential $Q$.
\end{corollary}

\begin{theorem}[$N=2$ topological-twist identities;
$\ClaimStatusProvedElsewhere$]
\label{thm:master-n2-topological-twist-topologisation}
Let $\mathfrak A$ be an N=2 superconformal vertex operator algebra at
central charge $c$, with generators
$(T, J, G^+, G^-)$ satisfying the standard N=2 SCA OPE
\cite[Chapter~5]{FrenkelBenZvi}.  Define the A-twist by
\[
T_A = T + \tfrac{1}{2}\partial J,
\qquad
Q_A = \oint \frac{dz}{2\pi i}\, G^+(z),
\qquad
G_A = \tfrac{1}{2}G^-.
\]
Then
\[
Q_A^2=0,
\qquad
\{Q_A,G_A(z)\}=T_A(z).
\]
The B-twist interchanges $+$ and $-$ and uses
$T_B=T-\frac12\partial J$.  These identities give T1 and make
holomorphic translation exact on BRST cohomology.  The remaining
entries of Convention~\ref{conv:master-topologisation} are properties
of the chosen completion and operadic chain model.
\end{theorem}

\begin{proof}
The defining N=2 OPEs include
\[
G^+(z) G^-(w) \;\sim\; \frac{2c/3}{(z-w)^3} + \frac{2 J(w)}{(z-w)^2}
 + \frac{2T(w) + \partial J(w)}{z - w} + \mathrm{reg}.
\]
Taking the zero-mode residue gives
\[
\{Q_A, G_A(z)\}
 = \tfrac{1}{2} \oint_z \frac{dw}{2\pi i}\, G^+(w) G^-(z)
 = \tfrac{1}{2}\bigl(2T(z) + \partial J(z)\bigr)
 = T_A(z).
\]
The OPE $G^+(z)G^+(w)$ is regular by the N=2 relations, hence
\[
Q_A^2 = \tfrac{1}{2}\{Q_A, Q_A\}
 = \tfrac{1}{2}\oint dz\, dw\, G^+(z) G^+(w) = 0
 \quad \text{for every } c.
\]
The B-twist follows from the OPE with the two supercurrents exchanged.
\end{proof}

\begin{remark}[The two topological differentials]
\label{rem:master-n2-twist-scope}
The A- and B-twists use the square-zero operators
$Q_A=G^+_0$ and $Q_B=G^-_0$.  Examples carrying this N=2 algebra
include
\begin{itemize}[noitemsep]
\item N=2 minimal models at $c = 3k/(k+2)$, $k \in \mathbb Z_{\geq 1}$
\textup{(}a discrete subfamily of class~$\mathsf M$\textup{);}
\item Kazama--Suzuki cosets $\widehat{\mathfrak g}_k /
\widehat{\mathfrak h}$ for Hermitian symmetric pairs;
\item N=2 sigma models on Calabi--Yau targets
\textup{(}the physical A- and B-model constructions\textup{).}
\end{itemize}
Each example inherits the identities of
Theorem~\ref{thm:master-n2-topological-twist-topologisation}.  Its
completed topological operad action is constructed in the specific
model.
\end{remark}

\begin{corollary}[Kazama--Suzuki cosets and the two twists;
$\ClaimStatusProvedElsewhere$]
\label{cor:master-n2-twist-kazama-suzuki}
Let $(G,H)$ be a Hermitian symmetric pair and let $k$ lie in the
admissible level range of Kazama--Suzuki.  Their coset
\[
M_{G/H}(k)
 =\operatorname{Com}\!\left(
 \widehat{\mathfrak h}_{k'},
 \widehat{\mathfrak g}_k\otimes
 \mathcal F^{\mathfrak g/\mathfrak h}
 \right)
\]
carries generators $(T,J,G^+,G^-)$ satisfying the N=2 OPE
\cite[equations~(3.27)--(3.30)]{KazamaSuzuki1989}.  Consequently its
A- and B-twists satisfy the two square-zero and stress-tensor
homotopy identities of
Theorem~\ref{thm:master-n2-topological-twist-topologisation}.
\end{corollary}

\begin{proof}
Apply the N=2 OPE calculation of
Theorem~\ref{thm:master-n2-topological-twist-topologisation} to the
Kazama--Suzuki generators.
\end{proof}

\begin{corollary}[Calabi--Yau sigma models and the A/B twists;
$\ClaimStatusConditional$]
\label{cor:master-n2-twist-cy-sigma}
Let the nonlinear sigma model of a smooth Calabi--Yau $d$-fold $Y$
admit an N$=(2,2)$ superconformal realization with chiral central
charge $c=3d$.  Each chiral half then carries the two differentials
of Theorem~\ref{thm:master-n2-topological-twist-topologisation}.  Their
physical realizations are Witten's A- and B-model twists
\cite{Witten91}.  A rigorous factorization model supplies the
completion, descent, and operadic action in T2--T5.
\end{corollary}

\begin{proof}
The formal implication follows from the N=2 OPEs.  The geometric
interpretation is the A/B-twist construction of~\cite{Witten91}.
\end{proof}

\begin{center}
\renewcommand{\arraystretch}{1.35}
\begin{tabular}{|p{3.0cm}|p{4.0cm}|p{6.5cm}|}
\hline
\textbf{Construction} & \textbf{Morphism} & \textbf{Statement} \\
\hline\hline
Universal reconstruction &
$\epsilon_\cA\colon\Omega B(\cA)\to\cA$ &
Theorem~A is the universal bar--cobar counit equivalence in the
complete enhanced associative Ran ambient under $H_1$. \\
\hline
Quadratic recognition &
$q_\cA\colon\cA^{\mathrm i}\to B(\cA)$ &
Theorem~B recognizes the locus on which $q_\cA$ is a
quasi-isomorphism for the chosen quadratic presentation; equivalently,
$\Omega(\cA^{\mathrm i})\to\cA$ is a quasi-isomorphism. \\
\hline
Fixed-$C$ second kind &
$D^{\mathrm{co}}(C\text{-}\mathrm{CoFact})\simeq
D^{\mathrm{ctr}}(C\text{-}\mathrm{ContraFact})$ &
For a fixed chiral CDG-coalgebra $C$, the co--contra equivalence is the
separate second-kind statement governed by its acyclic twisting
package. \\
\hline
Verdier--Koszul &
$\mathbb D_{\Ran}(q_\cA)\colon
\mathbb D_{\Ran}B(\cA)\to\cA^!$ &
Verdier duality sends the quadratic comparison
$q_\cA\colon\cA^{\mathrm i}\to B(\cA)$ to the displayed algebra map
in the finite-type dualizable ambient.  It is an equivalence under
$H_{\mathrm{CL}}\cup H_{\mathrm{VD}}$. \\
\hline
Chiral closed sector & $\cA\mapsto C^\bullet_{\mathrm{ch}}(\cA,\cA)$ &
Theorem~H concerns the chart cochain complex.  Comparison maps to
categorical Hochschild cochains, THH, or a physical bulk carry their
own source, target, and hypotheses. \\
\hline
\end{tabular}
\end{center}

\section{Scalar and collision normalizations}
\label{sec:master-concordance-normalizations}

The scalar conventions are those of
Chapter~\ref{ch:landscape-census}.  On the standard landscape,
\[
\begin{aligned}
\kappa(\mathcal H_\ell^{\oplus d}) &= d\ell,
&
\kappa(V_k(\fg)) &=
\frac{\dim(\fg)(k+h^\vee)}{2h^\vee},
&
\kappa(\mathrm{Vir}_c) &= \frac c2.
\end{aligned}
\]
On
$H_{\mathrm{diag}}^{g=1}+H_{W_N}^{\mathrm{DS/bar}}$, the principal
family has
\[
\kappa(\mathcal W_N)=c(H_N-1),
\qquad
H_N=\sum_{j=1}^{N}\frac1j.
\]
For the Virasoro shadow tower on the non-singular surface
$c(5c+22)\neq 0$,
\[
S_2=\frac c2,\qquad
S_3=2,\qquad
S_4=\frac{10}{c(5c+22)},\qquad
S_5=-\frac{48}{c^2(5c+22)},
\]
and
\[
\Delta=8\kappa S_4=\frac{40}{5c+22},\qquad
\Lambda={:}TT{:}-\frac{3}{10}\partial^2T,\qquad
\langle\Lambda|\Lambda\rangle=\frac{c(5c+22)}{10}.
\]
Thus $S_4$ is $10/[c(5c+22)]$; the number
$40/(5c+22)$ is the downstream discriminant $\Delta$.

The collision residue is written in trace-form normalization:
\[
r_{\mathcal H}^{\mathrm{coll}}(z)=\frac{k}{z},\qquad
r_{V_k(\fg)}^{\mathrm{coll}}(z)=\frac{k\Omega_{\mathrm{tr}}}{z},
\qquad
r_{\mathrm{Vir}_c}^{\mathrm{coll}}(z)=\frac{c/2}{z^3}+\frac{2T}{z}.
\]
For affine Kac--Moody, the KZ presentation
\[
r_{\mathrm{KZ}}(z)=\frac{\Omega}{(k+h^\vee)z},
\qquad
\Omega=2h^\vee\Omega_{\mathrm{tr}},
\]
is the monodromy normalization.  The trace-form collision function is
$k\Omega_{\mathrm{tr}}/z$.  Their values at $k=0$ are respectively
$\Omega/(h^\vee z)$ and~$0$.

\section{The seven-face atlas}
\label{sec:master-concordance-seven-faces}

\begin{theorem}[Seven-face atlas on the comparison locus;
\ClaimStatusConditional]
\label{thm:master-seven-face}
Let $\cA$ be a modular Koszul chiral algebra for which the following
comparison data are fixed:
\begin{enumerate}[label=\textup{(\alph*)}]
\item the bar-intrinsic MC element $\Theta_\cA=D_\cA-d_0$ of
Theorem~\ref{thm:mc2-bar-intrinsic};
\item the collision-residue/twisting comparison of
Theorem~\ref{thm:collision-residue-twisting};
\item the Vol~II line-operator comparison
\texttt{thm:dnp-bar-cobar-identification}, when the open-colour
$\mathsf{SC}^{\mathrm{ch,top}}$ hypotheses hold;
\item the Poisson-vertex comparison of
Theorem~\ref{thm:kz-classical-quantum-bridge}, in genus zero and
under its higher-genus hypotheses when those are invoked;
\item the GZ/Gaudin comparison of
Theorems~\ref{thm:gz26-commuting-differentials} and
\ref{thm:gaudin-yangian-identification};
\item the Drinfeld--Semenov-Tian-Shansky Yangian/Sklyanin
classical limit, with the affine scalar
$\hbar=(k+h^\vee)^{-1}$.
\end{enumerate}
The comparison locus means that the following target carriers and
comparison morphisms are defined:
\[
\begin{array}{rcl}
\mathcal T_1 &=& \mathrm{Tw}(\barBch(\cA),\cA),\\[2pt]
\mathcal T_2 &=&
 \End_{\mathrm{mer}}(\mathsf{Line}^{\mathrm{op}}_{\cA}),\\[2pt]
\mathcal T_3 &=&
 \mathcal R_{\mathrm{PVA}}\bigl(H^\bullet(\cA,Q)/(\hbar)\bigr),\\[2pt]
\mathcal T_4 &=&
 \End(\cV_{0,n}(\cA;V_1,\ldots,V_n))_{\mathrm{comm}},\\[2pt]
\mathcal T_5 &=&
 \mathcal R_{\mathrm{Dr}}\bigl(Y_\hbar(\fg)\bigr)
 \quad\text{on the affine Drinfeld row},\\[2pt]
\mathcal T_6 &=&
 \mathrm{Pois}_{\mathrm{Sk}}((\fg^!)^*)
 \quad\text{on the same classical row},\\[2pt]
\mathcal T_7 &=&
 \mathcal H_{\mathrm{Gaudin}}(\fg;z_1,\ldots,z_n).
\end{array}
\]
Here $\mathcal R_{\mathrm{Dr}}(Y_\hbar(\fg))$ is the classical
Drinfeld--Yangian carrier.  The Vol~III Hall--Drinfeld double and the
BKM denominator algebra are separate target objects.  The maps
$\rho_i:\mathcal T_1\to\mathcal T_i$ are precisely the comparison
morphisms supplied by the hypotheses above.
Then the seven expressions
\[
\begin{array}{ll}
\mathrm{F1}:&\text{universal twisting morphism,}\\
\mathrm{F2}:&\text{line-operator meromorphic tensor twist,}\\
\mathrm{F3}:&\text{Poisson-vertex classical }r\text{-matrix,}\\
\mathrm{F4}:&\text{commuting Hamiltonian seed,}\\
\mathrm{F5}:&\text{Yangian classical }r\text{-matrix,}\\
\mathrm{F6}:&\text{Sklyanin Poisson tensor,}\\
\mathrm{F7}:&\text{Gaudin Hamiltonian generator}
\end{array}
\]
are the images of the single class
\[
r_\cA(z)=\operatorname{Res}^{\mathrm{coll}}_{0,2}(\Theta_\cA)
\]
under the typed comparison morphisms $\rho_i$.  Equivalently,
\[
\bigl(\rho_1(r_\cA),\ldots,\rho_7(r_\cA)\bigr)
\in
\mathcal T_1\times_{\mathcal N_\cA}\cdots
\times_{\mathcal N_\cA}\mathcal T_7,
\]
where $\mathcal N_\cA$ is the normalization object generated by the
chosen rational function $r_\cA(z)$.  Thus the theorem identifies the
seven displayed images in the fibre product.  The carriers retain
their respective categories, and $\Omega B(\cA)$, $\cA^!$, and
$Z^{\mathrm{der}}_{\mathrm{ch}}(\cA)$ retain their typed definitions.
\end{theorem}

\begin{proof}
Theorem~\ref{thm:collision-residue-twisting} identifies
$\operatorname{Res}^{\mathrm{coll}}_{0,2}(\Theta_\cA)$ with the
universal twisting morphism
$\pi_\cA\colon B(\cA)\to\cA$, giving F1.  The Vol~II line-operator
theorem sends the same twisting datum to the meromorphic tensor
product of open-colour line operators, giving F2 on its hypotheses.
The Poisson-vertex bridge reads the genus-zero bracket coefficients
as the same collision residue, giving F3.  The GZ theorem inserts
the depth-$k$ collision residues into the sphere Hamiltonians, and
the Gaudin theorem identifies the affine depth-one specialization
with the usual Gaudin generator; these are F4 and F7.  On the affine
Drinfeld row, Drinfeld's Yangian theorem and the
Semenov-Tian-Shansky Sklyanin bracket place the classical Yangian
$r$-matrix and the Sklyanin tensor over the same rational
normalization after the scalar $\hbar=(k+h^\vee)^{-1}$ is fixed; these
are F5 and F6.  The Vol~III Hall--Drinfeld/CoHA and BKM comparisons
require their own morphisms.  Each arrow is therefore a morphism out of the same
pair
$(\Theta_\cA,\operatorname{Res}^{\mathrm{coll}}_{0,2})$, so the
seven images agree in the stated fibre product.
\end{proof}

\section{Face anchors}
\label{sec:master-concordance-face-locations}

\begin{center}
\renewcommand{\arraystretch}{1.35}
\begin{tabular}{|p{2.8cm}|p{3.5cm}|p{3.5cm}|p{3.5cm}|}
\hline
\textbf{Face} & \textbf{Vol~I anchor} & \textbf{Vol~II anchor}
& \textbf{Vol~III anchor} \\
\hline\hline
F1: twisting morphism &
\texttt{thm:mc2-bar-intrinsic};
\texttt{thm:collision-residue-twisting} &
\texttt{line-operators.tex};
open-colour bar--cobar &
Stage-$2$ $E_1$-chiral shadow of
$\Phi_d^{(\Sigma,C)}$ \\
\hline
F2: line-operator twist &
\texttt{thm:collision-residue-twisting} &
\texttt{thm:dnp-bar-cobar-identification} &
line/boundary specialization after $\Sp^{\mathrm{ch}}_{\Sigma,C}$ \\
\hline
F3: PVA $r$-matrix &
\texttt{thm:kz-classical-quantum-bridge} &
\texttt{pva\_descent.tex} &
\texttt{cy\_to\_chiral.tex}: genus-zero shadow \\
\hline
F4: commuting Hamiltonians &
\texttt{thm:gz26-commuting-differentials} &
HT sphere/interface Hamiltonians &
CY shadow Hamiltonians on constructed loci \\
\hline
F5: Drinfeld Yangian $r$-matrix &
\texttt{yangians\_drinfeld\_kohno.tex} &
dg-shifted factorization Yangian &
Hall--Drinfeld/CoHA positive half, a target distinct from
$Y_\hbar(\mathfrak g_{\Delta_5})$ \\
\hline
F6: Sklyanin tensor &
\texttt{thm:yangian-sklyanin-quantization} &
spectral-braiding core &
toroidal/elliptic Hall Sklyanin form \\
\hline
F7: Gaudin generator &
\texttt{thm:gaudin-yangian-identification} &
closed--open line Hamiltonians &
CY holographic datum on the Stage-$2$ curve \\
\hline
\end{tabular}
\end{center}

\section{Face bridge theorems}
\label{sec:master-concordance-new-theorems}

\begin{center}
\renewcommand{\arraystretch}{1.35}
\begin{tabular}{|p{5.2cm}|p{2.0cm}|p{6.4cm}|}
\hline
\textbf{Label} & \textbf{Volume} & \textbf{Mathematical content} \\
\hline\hline
\texttt{thm:dnp-bar-cobar-identification} & II &
The open-colour line-operator package is the bar--cobar twisting
package on the chirally Koszul locus. \\
\hline
Theorem~\ref{thm:gz26-commuting-differentials} & I &
The genus-zero shadow of $\Theta_\cA$ gives commuting Hamiltonians;
for $\mathcal W_N$ the existence is algebraic and the term-by-term
GZ comparison is a separate normalization problem. \\
\hline
Theorem~\ref{thm:kz-classical-quantum-bridge} & I &
The PVA $\lambda$-bracket determines the genus-zero bar differential
and its collision residue; higher-genus geometric comparison uses the
stated BV/BRST hypotheses. \\
\hline
Theorem~\ref{thm:gaudin-yangian-identification} & I &
On affine Kac--Moody, the depth-one Hamiltonian is
$(k+h^\vee)^{-1}$ times the Gaudin Hamiltonian; higher depths are
transferred $A_\infty$ operations. \\
\hline
Theorem~\ref{thm:yangian-sklyanin-quantization} & I &
The Drinfeld/Sklyanin classical limit is external; the local
contribution is the scalar identification
$\hbar=(k+h^\vee)^{-1}$ across the affine comparison surfaces. \\
\hline
Theorem~\ref{thm:shadow-depth-operator-order} & I &
The GZ operator order is controlled by collision depth
$k_{\max}=p_{\max}-1$, while the shadow class is controlled by
$r_{\max}$. \\
\hline
\end{tabular}
\end{center}

\section{The five named theorems}
\label{sec:master-concordance-theorem-surfaces}

\begin{center}
\renewcommand{\arraystretch}{1.35}
\begin{tabular}{|p{1.1cm}|p{4.0cm}|p{8.4cm}|}
\hline
\textbf{Thm} & \textbf{Name} & \textbf{Typed statement} \\
\hline\hline
A & Universal bar--cobar reconstruction &
Vol~I constructs the chiral bar--cobar adjunction.  Its unit and the
counit
$\epsilon_\cA\colon\Omega B(\cA)\to\cA$ are equivalences in the
complete enhanced associative Ran ambient under $H_1$, with the
counit carrying target $\cA$.  Verdier duality separately forms the
algebra $\mathbb D_{\Ran}B(\cA)$.  Vol~II carries the open-colour
specialization in $\mathsf{SC}^{\mathrm{ch,top}}$.  Vol~III carries the
Stage-$2$ specialization to a reference curve:
$\Phi_d^{(\Sigma_{d-1},C)}
=\mathrm{Sp}^{\mathrm{ch}}_{\Sigma_{d-1},C}\circ\Phi_d^{\mathrm{FA}}$. \\
\hline
B & Quadratic Koszul recognition &
For a chosen quadratic presentation
$\cA=T(V)/(R)$, Theorem~B recognizes the locus on which
$q_\cA\colon\cA^{\mathrm i}\to B(\cA)$ is a quasi-isomorphism.
Equivalently, $\Omega(\cA^{\mathrm i})\to\cA$ is a quasi-isomorphism;
the chiral transport uses the package $H_{\mathrm{CL}}$. \\
\hline
-- & Fixed-$C$ second-kind equivalence &
For a fixed chiral CDG-coalgebra $C$, the equivalence
$D^{\mathrm{co}}(C\text{-}\mathrm{CoFact})\simeq
D^{\mathrm{ctr}}(C\text{-}\mathrm{ContraFact})$
is the independent second-kind equivalence governed by
$\mathsf{Tw}^{\mathrm{ch}}_{\mathrm{acyc}}(C,A,\tau)$. \\
\hline
C & Centre-local-system complementarity &
C0 identifies the strict-flat degree-zero fibre with
$\mathcal Z(\cA)$; C1 decomposes
$\mathbf C_g(\cA)=R\Gamma(\overline{\mathcal M}_g,\mathcal Z(\cA))$
by the represented Verdier involution.  A supplied
$\iota_Z^{\mathrm{der}}$, followed by~$H^0$ and C0, connects the
Theorem~H derived centre to this coefficient system.  Theorem~D gives
the scalar trace
$K^\kappa(\cA)=\kappa(\cA)+\kappa(\cA^!)$. \\
\hline
D & Deformation, Hodge, and graph traces &
$\operatorname{Obs}^{\mathrm{def}}_g(\cA)$ lies in
$H^2(\Def_g(\cA))$ for $g\geq2$.  Under~$H_D^1$, the pointed class
$\operatorname{tr}_1\operatorname{Obs}^{\mathrm{def}}_{1,1}
=\kappa(\cA)\lambda_1$.  Under~$H_D^K$,
$\mathfrak O_g^K=\kappa(\cA)\lambda_{-1}(\mathbb E_g)$ and
$\operatorname{ch}_g(\mathfrak O_g^K)=(-1)^g\kappa(\cA)\lambda_g$.
The package~$H_D^{\mathrm{tr}}$ supplies their derived comparison.
Under~$H_D^{\mathrm{graph}}$,
$F_g(\cA)=F_g^{\mathrm{sc}}(\cA)
+\delta F_g^{\mathrm{cross}}(\cA)$, where
$F_g^{\mathrm{sc}}(\cA)=\kappa(\cA)\lambda_g^{\mathrm{FP}}$. \\
\hline
H & Family-indexed chiral Hochschild support &
For a finite set $S\subset\mathbb Z$, the package $H_H(\cA;S)$
constructs an incidence-compatible strong deformation retract from a
complete chart-cochain model to a complex supported in $S$; it then
gives
$\operatorname{Supp}\ChirHoch^\bullet(\cA)\subseteq S$.
The bounded vertex-algebra benchmarks of
Bakalov--De Sole--Kac have supports $\{0,2,3\}$ for Virasoro and
dimension vector $(2,1,0,\ldots)$ for the rank-one even free
superboson.  Chartwise transport uses a separately named comparison
map. \\
\hline
\end{tabular}
\end{center}

\section{The shadow classes}
\label{sec:master-concordance-classification}

For the four displayed census families, the pole order and shadow
depth are the following separate invariants.
\[
\begin{array}{c|c|c|c|l}
\textbf{Class} & p_{\max} & k_{\max} & r_{\max} &
\textbf{Witnesses}\\
\hline
G & 2 & 1 & 2 & \mathcal H_k,\; V_\Lambda,\; \text{free fermion}\\
L & 2 & 1 & 3 & V_k(\fg)\text{ at non-critical level}\\
C & 1 & 0 & 4 & \beta\gamma_\lambda,\; bc_\lambda
\text{ in the standard conformal-weight family}\\
M & 4\text{ or }2N & 3\text{ or }2N-1 & \infty &
\mathrm{Vir}_c,\; \mathcal W_N,\; \mathrm{BP}_k
\end{array}
\]
Thus $k_{\max}=p_{\max}-1$ records the degree-two collision residue,
while $r_{\max}$ records the full obstruction tower.  The
$\beta\gamma$ row is the separating witness:
\[
p_{\max}(\beta\gamma)=1,\qquad
k_{\max}(\beta\gamma)=0,\qquad
r_{\max}(\beta\gamma)=4.
\]
The class-$C$ quartic obstruction is carried by the standard
conformal-weight family, with
$S_2=6\lambda^2-6\lambda+1$, $S_3=0$,
$S_4=-5/12$, and $S_r=0$ for $r\geq5$.
The scope of this table is exactly the listed witnesses.  The
class-$\mathsf B$ Hall/Borcherds depth uses the distinct filtration
defined in the K3 row of Chapter~\ref{ch:landscape-census}.

\section{The three depths}
\label{sec:master-concordance-three-invariants}

For a chosen minimal strong generating set $\{a_\alpha\}$,
\[
p_{\max}(\cA)
=\max_{\alpha,\beta}
\{n\mid a_\alpha(z)a_\beta(w)
\text{ has a nonzero }(z-w)^{-n}\text{ term}\}.
\]
The collision depth is
\[
k_{\max}(\cA)
=\max\{k\mid
\operatorname{Res}^{\mathrm{coll}}_{0,2}(\Theta_\cA)
\text{ has a nonzero }z^{-k}\text{ term}\},
\]
and the bar propagator absorbs one pole:
\[
k_{\max}(\cA)=p_{\max}(\cA)-1.
\]
The shadow depth
\[
r_{\max}(\cA)=
\max\{r\geq2\mid
\Theta_\cA^{\leq r}
\text{ has a nonzero degree-}r\text{ projection}\}
\in\{2,3,4,\infty\}
\]
depends on higher composites and normal-ordered contact terms.  Hence
the GZ Hamiltonian operator order is controlled by $k_{\max}$, while
the quadrichotomy $G/L/C/M$ is controlled by $r_{\max}$.

\section{Complementarity scalars}
\label{sec:master-concordance-complementarity-scalars}

The scalar Verdier sum
\[
K^\kappa(\cA):=\kappa(\cA)+\kappa(\cA^!)
\]
is family-dependent. On the Theorem~C surface it is the scalar trace
shadow of the two \(\mathrm{C}_1\) Verdier eigensummands of
$\mathbf C_g(\cA)=R\Gamma(\overline{\mathcal M}_g,\mathcal Z(\cA))$
after the Theorem~D modular trace
\textup{(}Proposition~\ref{prop:scalar-conductor-c1-trace-shadow}\textup{)};
the number records the final scalar image.  The Theorem~H derived
centre reaches the C0 coefficient system through the supplied brace
comparison~$\iota_Z^{\mathrm{der}}$, its $H^0$ shadow, and the strict-flat
fibre--centre identification.  The trace map and the \(\mathrm{C}_2\)
shifted-symplectic datum remain preceding typed objects. Here $\cA^!$ is the
\emph{algebra-level Verdier dual}: Feigin--Frenkel level shift
$k\mapsto -k-2h^\vee$ for
class~$\mathsf L$, $c\mapsto 26-c$ for the Virasoro row of
class~$\mathsf M$, and ghost-pair flip for class~$\mathsf C$.
For the Bershadsky--Polyakov row, write
$H_{\mathrm{BP}}^{\mathrm{DS/bar}}$ for the package consisting of a
regular level $k\ne-3$, a finite-type bar dual or a completed
continuous nondegenerate Verdier pairing, and subregular hook-type
DS/bar transport.  On this package the comparison theorem identifies
$\mathrm{BP}_k^!\simeq\mathrm{BP}_{-k-6}$.
Convention~\ref{conv:A-two-kappa-shriek} introduces a separate
chain-level scalar $\nu_A$ through
\[
 d_Ah_A+h_Ad_A
 =\nu_A(\mathrm{id}-\iota_Ap_A).
\]
An identity $\nu_A=K^\kappa(A)$ requires a chain map from this
deformation retract to the scalar trace complex.  The current family
values
\[
 \nu_{\mathsf G}=1,
 \quad
 \nu_{\mathsf L}=2(k+h^\vee),
 \quad
 \nu_{\mathsf C}=1,
 \quad
 \nu_{\mathsf M}=\frac{c(5c+22)}{10},
 \quad
 \nu_{\mathsf B}=8
\]
are conditional construction targets in
Remark~\ref{rem:A-h-three-paths}.  The scalar Verdier sums computed
from the displayed partner formulas are listed separately below.
\begin{center}
\renewcommand{\arraystretch}{1.35}
\begin{tabular}{|p{3.6cm}|p{3.2cm}|p{3.2cm}|p{3.0cm}|}
\hline
\textbf{Family} & \textbf{Verdier partner} &
\textbf{$K^\kappa$} & \textbf{Source formula} \\
\hline\hline
Heisenberg $\mathcal H_k$ &
curved $\mathrm{Sym}^{\mathrm{ch}}(V^*)$ &
$0$ & $\kappa^!=-k$ \\
\hline
Affine $V_k(\fg)$ &
Feigin--Frenkel level $-k-2h^\vee$ &
$0$ &
$\kappa(V_k(\fg))=\dim(\fg)(k+h^\vee)/(2h^\vee)$ \\
\hline
$\beta\gamma_\lambda$ &
$bc_\lambda$ &
$0$ &
$c_{bc}=-c_{\beta\gamma}$ and $\kappa=c/2$ \\
\hline
Virasoro $\mathrm{Vir}_c$ &
$\mathrm{Vir}_{26-c}$ &
$13$ &
$\kappa=c/2$ \\
\hline
Principal $\mathcal W_3^k$ &
$\mathcal W_3^{-k-6}$ on $H_{W_3}^{\mathrm{DS/bar}}$ &
$250/3$ on
$H_{\mathrm{diag}}^{g=1}+H_{W_3}^{\mathrm{DS/bar}}$ &
$K_3^c=100$ exactly;
$\kappa=(5/6)c$ on $H_{\mathrm{diag}}^{g=1}$ \\
\hline
Bershadsky--Polyakov $\mathrm{BP}_k$ &
$\mathrm{BP}_{-k-6}$ on $H_{\mathrm{BP}}^{\mathrm{DS/bar}}$ &
open &
$K^c=50$ exactly; genus-one $\kappa$ calculation open \\
\hline
Mukai-enhanced K3 Heisenberg &
Vol~III $\mathsf B$-row witness &
$8$ &
$2c_+(\mathrm{Mukai}(K3))=8$ exactly;
chiral interpretation conjectural under $H_{\mathsf B}$ \\
\hline
\end{tabular}
\end{center}
The unconditional scalar values in the displayed Vol.~I witness rows
are
\[
\{0,\,13\}.
\]
The exact principal central conductor is $K_3^c=100$.  On
$H_{\mathrm{diag}}^{g=1}+H_{W_3}^{\mathrm{DS/bar}}$, its modular
image is the conditional value $K_3^\kappa=250/3$, giving the
comparison set
\[
\{0,\,13,\,250/3\}.
\]
The BP value $25/3$ is conditional on its genus-one curvature
calculation, and the Mukai-enhanced value~$8$ is conditional on
$H_{\mathsf B}$.

\begin{remark}[Bershadsky--Polyakov conformal-vector conventions]
\label{rem:master-bp-convention-separation}
The Fehily--Kawasetsu--Ridout conformal vector gives the rational
identities
\[
c_{\mathrm{BP}}(k)
=-\frac{(2k+3)(3k+1)}{k+3},
\qquad
c_{\mathrm{BP}}(k)+c_{\mathrm{BP}}(-k-6)=50,
\]
Fehily--Kawasetsu--Ridout define an ordinary bosonic vertex algebra,
so $J,T,G^+,G^-$ are even.  Their reciprocal-weight diagnostic is
$1+2/3+2/3+1/2=17/6$.  The modular characteristic is
\ClaimStatusOpen{} pending the complete genus-one minimal-DS curvature
calculation.  A result $\kappa_{\mathrm{BP}}=c_{\mathrm{BP}}/6$
would imply $K^\kappa_{\mathrm{BP}}=25/3$.
The explicitly shifted secondary conformal vector carries the
computed scalar identity
\[
c_{\mathrm{BP}}^{\mathrm{shift}}(k)
=2-\frac{24(k+1)^2}{k+3},
\qquad
c_{\mathrm{BP}}^{\mathrm{shift}}(k)
+c_{\mathrm{BP}}^{\mathrm{shift}}(-k-6)=196.
\]
The value~$196$ belongs to this computed-secondary scalar convention,
whose role is convention comparison.  Both rational functions have a
common pole at $k=-3$, the fixed point of $k\mapsto-k-6$.
The algebra-level same-family comparison uses the first conformal
vector and the conditional package $H_{\mathrm{BP}}^{\mathrm{DS/bar}}$.
\end{remark}

The exact Mukai-lattice number is
$2c_+(\mathrm{Mukai}(K3))=8$; its proposed chiral interpretation is
$K^{\kappa_{\mathrm{ch}}}=8$ under $H_{\mathsf B}$.  The Borcherds
weight $\kappa_{\mathrm{BKM}}(\Delta_5)=5$ is a distinct scalar
invariant of the automorphic product.
Self-dual lattice VOAs have their own lattice-dual scalar, e.g.
$V_{\Lambda_{24}}^!\cong V_{\Lambda_{24}}$ gives
$\kappa+\kappa^!=48$; this is the lattice-dual scalar rather than a
Feigin--Frenkel level-shift complementarity constant.

\begin{remark}[Parameter involutions and KSDual candidates;
scalar relations $\ClaimStatusComputed$, object data
$\ClaimStatusConditional$]
\label{rem:ksdual-fixed-points-bucket}
\index{KSDual sublocus!archetype-by-archetype fixed point}%
Each row carries a scalar parameter involution or orbit relation.  The
table records its fixed parameters.  An object of $\KSDual$ is the
triple $(A,\varphi,H)$ of Definition~\ref{def:mr-ksdual}; construction
of $\varphi\colon A\simeq A^!$ and its coherence $H$ is the
object-level problem.
\begin{center}
\renewcommand{\arraystretch}{1.25}
\begin{tabular}{@{}llll@{}}
\toprule
Archetype & Family & Parameter relation & Fixed-parameter candidate\\
\midrule
$\mathsf G$ & $\cH_k$ & $k\mapsto -k$ & $k = 0$\\
$\mathsf L$ & $V_k(\fg)$ & $k\mapsto -2h^\vee - k$ & $k = -h^\vee$ \textup{(}critical\textup{)}\\
$\mathsf C$ & $\beta\gamma_\lambda \leftrightarrow bc_\lambda$ & family exchange &
an equivalence $\beta\gamma_\lambda\simeq bc_\lambda$\\
$\mathsf M_{\mathrm{Vir}}$ & $\Vir_c$ & $c\mapsto 26 - c$ & $c = 13$\\
$\mathsf M_{\cW_3}$ & $\cW_3^k$ & $c\mapsto 100 - c$ & $c = 50$\\
$\mathsf M_{\BP}$ & $\BP_k$ & $k\mapsto -k - 6$ &
$k=-3$ \,\textup{(}fixed pole; comparison boundary\textup{)}\\
$\mathsf B$ & Mukai-K3 & lattice involution & K3 Mukai-self-dual\\
\bottomrule
\end{tabular}
\end{center}
The scalar companion midpoints of the three displayed
$\mathsf M$-families are
\[
c^{\mathrm{mid}}_{\mathrm{Vir}}=13=26/2,
\qquad
c^{\mathrm{mid}}_{\cW_3}=50=100/2,
\qquad
c^{\mathrm{mid}}_{\mathrm{BP}}=25=50/2.
\]
At an object-level KSDual point, the map
$\varphi\colon A\xrightarrow{\sim}A^!$ is the level-$2$ equivalence.
An equivariant modular realization transports this involution to the
level-$5$ trace data.
For BP, $H_{\mathrm{BP}}^{\mathrm{DS/bar}}$ identifies the regular
companion pairs at level~$2$; the fixed parameter $k=-3$ lies at the
common pole and marks the boundary of this regular comparison locus.

At $k=0$, $\cH_0$ is the commutative rank-one current vertex algebra;
its collision form vanishes.  At the affine fixed parameter
$k=-h^\vee$, the Feigin--Frenkel centre
$Z(V_{-h^\vee}(\fg))$ is infinite-dimensional and the candidate
rescaling scalar $\nu_{\mathsf L}=2(k+h^\vee)$ vanishes.  A
deformation retract using $\nu_{\mathsf L}^{-1}$ is therefore defined
on the localized noncritical parameter ring.  Theorem~H at the
critical point requires its own family support package.

\emph{Complex scalar-midpoint fibres for $\cW_3$ and $\BP$.}
By the Fateev--Lukyanov principal-DS formula of
Remark~\ref{rem:koszul-conductor-explicit},
\[
c_{\cW_3}(k)
=2-\frac{24(k+2)^2}{k+3}
=50-24\left((k+3)+\frac{1}{k+3}\right).
\]
Consequently
\[
c_{\cW_3}(k)=50
\qquad\Longleftrightarrow\qquad
(k+3)^2+1=0
\qquad\Longleftrightarrow\qquad
k=-3\pm i.
\]
These two regular principal-DS levels are exchanged by
$k\mapsto-k-6$.  The standard BP scalar midpoint is
$c^{\mathrm{mid}}_{\mathrm{BP}}=25$.
Indeed,
\[
-\frac{(2k+3)(3k+1)}{k+3}=25
\qquad\Longleftrightarrow\qquad
k=-3\pm2i.
\]
These two regular BP levels are exchanged by $k\mapsto-k-6$.
The parameter involution fixes $k=-3$, where the central-charge
function has its pole.  Thus $c=25$ is the regular scalar-midpoint
fibre, and $H_{\mathrm{BP}}^{\mathrm{DS/bar}}$ supplies its
same-family Verdier--Koszul interpretation.
\end{remark}

\section{Cross-volume specialization}
\label{sec:master-concordance-cross-volume}

Vol~I contains the bar--cobar reconstruction arrows.  Vol~II and
Vol~III carry specialization functors from the reconstructed data.

\begin{enumerate}[label=\textup{(\roman*)}]
\item Vol~II uses the two-coloured dioperad
$\mathsf{SC}^{\mathrm{ch,top}}$.  Its closed-output operations have
closed inputs, while its open-output operations admit both closed and
open inputs.  The open-colour bar--cobar comparison supplies line
modules over the Koszul dual, and the closed cochain object is the
chiral Hochschild complex.
\item Vol~III uses the two-stage factorisation
\[
\Phi_d^{(\Sigma_{d-1},C)}
=\mathrm{Sp}^{\mathrm{ch}}_{\Sigma_{d-1},C}
\circ\Phi_d^{\mathrm{FA}}.
\]
Stage~$1$ produces an $E_d$ holomorphic factorisation algebra.
Stage~$2$ integrates along $\Sigma_{d-1}$ and restricts to a reference
curve $C$, producing an $E_1$-chiral object.  Its type is a CY-to-open
specialization functor.
\item The shared Theorem~A surface is the curve-side $E_1$ chiral
shadow after Stage~$2$.  A single CY$_d$ category may therefore have
many $E_1$ shadows, indexed by $(\Sigma_{d-1},C)$.
\end{enumerate}

The six-level Beilinson fork of
Convention~\ref{conv:master-primitive-datum} carries a typed vertical
comparison morphism at every level to its Vol~II or Vol~III companion
\textup{(}\texttt{thm:mr-vertical-dictionary} in
\texttt{master\_reconstruction.tex}\textup{).}  Each equivalence claim
uses its displayed package $H_{Vj}$.  The dedicated chapters
exist for levels $0$ and $4$
\textup{(}\texttt{vertical\_equivalence\_level\_0.tex},
\texttt{vertical\_equivalence\_level\_4.tex}\textup{).}  The
intermediate levels $1$ and $2$ admit the following lemma-level
inscriptions, stating the chart-algebra and bar-coalgebra vertical
equivalences as direct consequences of the dictionary theorem.

\begin{lemma}[Level-$1$ vertical equivalence: chart algebra
$A_b$ $\leftrightarrow$ $\Phi_d^{\mathrm{FA}}$ at chosen vacuum;
$\ClaimStatusConditional$ on hypothesis package $H_{V1}$]
\label{lem:master-vertical-level-1}
Suppose the Open factorization compact generator $b\in\mathsf C^\omega$
over $(X,D,\tau)$ corresponds via the vacuum-target
correspondence to a CY-$d$ target $c$ in the CY-side category
$\mathrm{CY}_d\text{-Cat}$. Then the chart algebra
\[
A_b \;=\; \RHom_{\mathsf C}(b,b)
\quad\text{(Open, level~$1$)}
\]
is quasi-isomorphic as an $E_1$-chiral algebra to the evaluation of
the Vol~III Stage-$1$ factorisation-algebra functor at $c$:
\[
A_b \;\simeq\; \Phi_d^{\mathrm{FA}}(\cC)\bigl|_c
\quad\text{(Vol~III, level~$1$).}
\]
The equivalence is conditional on hypothesis package $H_{V1}$ of
\texttt{thm:mr-vertical-dictionary}: Kontsevich--Tamarkin
$E_d$-formality of $\mathrm{CY}_d\text{-Cat}$ for $d \leq 3$, via
three-vanishing of the obstruction class.
\end{lemma}

\begin{proof}
The vacuum-target correspondence supplies a chosen identification
$b \leftrightarrow c$ at chart level. By the
$E_d$-formality~\cite{Kon99,CalaqueWillwacher21} of the CY-side
category, $\Phi_d^{\mathrm{FA}}(\cC)\bigl|_c$ is the $E_d$-formal
factorisation-algebra evaluation at $c$, equivalent to the
endomorphism algebra
$\End_{\Phi_d^{\mathrm{FA}}(\cC)}(c)$ on the CY side. This equals
$\End_{\cC^{\mathrm{op}}}(b)$ on the Open side by the chosen
identification, hence equals $A_b$.
\end{proof}

\begin{lemma}[Level-$2$ vertical equivalence: bar coalgebra
$B(A_b)$ $\leftrightarrow$ $\mathrm{Sp}^{\mathrm{ch}}_{\Sigma_{d-1},
C}$ specialisation; $\ClaimStatusConditional$ on
$H_{V1} \cup H_{V2}$]
\label{lem:master-vertical-level-2}
Under the vacuum-target correspondence of
Lemma~\ref{lem:master-vertical-level-1}, the Open bar coalgebra at
level~$2$ and the Vol~III two-stage specialisation at level~$2$
agree:
\[
B(A_b)
\;\simeq\;
B\bigl(\mathrm{Sp}^{\mathrm{ch}}_{\Sigma_{d-1}, C}
       \circ \Phi_d^{\mathrm{FA}}(\cC)\bigr)
\quad\text{as }E_1\text{-chiral coalgebras over }(\mathrm{Ass}^{\mathrm{ch}})^!.
\]
The equivalence is conditional on $H_{V1}$
\textup{(}Lemma~\ref{lem:master-vertical-level-1}\textup{)} and
$H_{V2}$ \textup{(}well-definedness of the Stage-$2$ specialisation
functor $\mathrm{Sp}^{\mathrm{ch}}_{\Sigma_{d-1}, C}$ for the chosen
$(\Sigma_{d-1}, C)$\textup{)}, both of
\texttt{thm:mr-vertical-dictionary}.
\end{lemma}

\begin{proof}
By Lemma~\ref{lem:master-vertical-level-1},
$A_b \simeq \Phi_d^{\mathrm{FA}}(\cC)\bigl|_c$. The Stage-$2$
specialisation $\mathrm{Sp}^{\mathrm{ch}}_{\Sigma_{d-1}, C}$ at the
chosen reference curve $C$ produces an $E_1$-chiral algebra on $C$
whose evaluation at the chosen target $c$ agrees with the Stage-$1$
evaluation, i.e.,
$\mathrm{Sp}^{\mathrm{ch}}_{\Sigma_{d-1}, C}(\Phi_d^{\mathrm{FA}}(\cC))\bigl|_c
\simeq A_b$ as $E_1$-chiral algebras
\textup{(}hypothesis $H_{V2}$\textup{).} Apply the bar functor
$B = T^c(s^{-1} \overline{(\cdot)})$ to both sides: by functoriality,
$B(A_b) \simeq B(\mathrm{Sp}^{\mathrm{ch}}_{\Sigma_{d-1}, C}
\circ \Phi_d^{\mathrm{FA}}(\cC))$ as $E_1$-chiral coalgebras over
$(\mathrm{Ass}^{\mathrm{ch}})^!$.
\end{proof}

\begin{lemma}[Level-$5$ vertical equivalence: modular trace
$\leftrightarrow$ universal Borcherds weight;
$\ClaimStatusConditional$ on $H_{V5}$]
\label{lem:master-vertical-level-5}
At level $5$ of the Beilinson tower, the Open-side modular trace
$\mathrm{Tr}_{\cC^{\mathrm{op}}}(A_b)$ on the modular shadow
$\overline{\cM}_{g, n}$ and the Vol~III universal Borcherds weight
$\kappa_{\mathrm{BKM}}(\Phi_N) = c_N(0)/2$ are joint scalar
invariants of the same modular datum. For target
$X = K3 \times E$ with class~$\mathsf B$ chart algebra
$A_b = $ Mukai-K3-Heisenberg, both project to the Borcherds product
$\Delta_5$:
\begin{enumerate}[label=\textup{(\arabic*)}]
\item the Open modular trace has scalar trace equal to the BPS counting
modular form
$Z_{\mathrm{BPS}}^{K3 \times E} =
(\Phi_{10}^{\mathrm{un}})^{-1} = \Delta_5^{-2}$ as an element of
$M_{-10}^{\mathrm{mero}}(\mathrm{Sp}_4(\mathbb Z))$
\textup{(}meromorphic modular form of weight $-10$ on
the Siegel paramodular group $K(1)$\textup{);}
\item the Vol~III universal Borcherds weight reads
$\kappa_{\mathrm{BKM}}(\Delta_5) = c_5(0)/2 = 5$ as the weight of
$\Delta_5$ in $S_5(K(1))$
\textup{(}weight-$5$ cusp form\textup{);}
\end{enumerate}
both invariants are jointly determined by $\Delta_5$, with~\textup{(1)}
the inverse partition function and~\textup{(2)} the modular weight.
\end{lemma}

\begin{proof}
The Borcherds product $\Delta_5$ for the K3-BKM algebra
$\mathfrak g_{\Delta_5}$ is the additive lift
$\mathrm{Grit}(\eta^9 \vartheta_1) \in S_5(K(1))$
\textup{(}\cite[Thm.~6.1]{Gritsenko1999}\textup{).} The Borcherds
denominator identity reads
\[
e^{2\pi i \langle \rho, Z\rangle}\,
\prod_{\alpha \in \Phi^+_{\Delta_5}}
\bigl(1 - e^{2\pi i \langle \alpha, Z\rangle}\bigr)^{\mathrm{mult}(\alpha)}
\;=\;\Delta_5(Z),
\]
exhibiting $\Delta_5$ as the Weyl--Kac denominator of
$\mathfrak g_{\Delta_5}$.

\textup{(}1\textup{)} The Open-side modular trace on the K3$\times$E
modular shadow $\overline{\cM}_{1, 1}^{K3\times E}$ produces, via
clutching on the open category, the BPS partition function which
equals the inverse of the squared Borcherds form
$Z_{\mathrm{BPS}}^{K3 \times E} = (\Phi_{10}^{\mathrm{un}})^{-1}$
\textup{(}Dijkgraaf--Verlinde--Verlinde 1997 1/4-BPS counting on
K3$\times T^2$\textup{);} this equals $\Delta_5^{-2}$ by the
identification $\Phi_{10}^{\mathrm{un}} = \Delta_5^2$
\textup{(}Gritsenko 1999\textup{).}

\textup{(}2\textup{)} The Vol~III universal Borcherds weight identity
$\kappa_{\mathrm{BKM}}(\Phi_N) = c_N(0)/2$ for $N \in \{1, 2, 3, 4, 6\}$
reads off the modular weight of $\Phi_N$ from the constant Fourier
coefficient of the input weak Jacobi form. For
$\Delta_5$, the additive construction uses
$\mathrm{Grit}(\eta^9\vartheta_1)$, whereas the multiplicative
Borcherds product uses the weak Jacobi form $\phi_{0,1}$ with
constant coefficient $f_1(0,0)=10$.  The multiplicative weight formula
therefore gives
$\kappa_{\mathrm{BKM}}(\Delta_5)=10/2=5$
\cite[Example~2.4, equation~(2.16)]{GritsenkoNikulin1998}.

Both \textup{(1)} and \textup{(2)} are scalar invariants determined
by the single modular form $\Delta_5$: the first via inversion and
squaring, the second via reading off the modular weight. The vertical
equivalence at level~$5$ identifies these two projections as the
Open- and CY-side scalar shadows of the same Borcherds datum, modulo
the Maloney--Witten + Manschot--Moore Farey-tail gravity-line
completion of hypothesis $H_{V5}$
\textup{(}\texttt{thm:mr-vertical-dictionary}\textup{).}
\end{proof}

\section{Boundary of the dictionary}
\label{sec:master-concordance-boundary}

The seven-face atlas is a theorem on the comparison locus of
Theorem~\ref{thm:master-seven-face}.  Its extension requires five
further constructions.
\begin{enumerate}
\item The term-by-term comparison between the $\mathcal W_N$
Hamiltonians produced by $\Theta_{\mathcal W_N}$ and the operators
of Gaiotto--Zeng is a normalization problem beyond the algebraic
commutativity theorem.
\item Genus $g\geq2$ multi-weight algebras carry the cross-channel
correction
$\delta F_g^{\mathrm{cross}}(\cA)$; the seven-face scalar statement
uses the uniform-weight projection.
\item Arbitrary-nilpotent non-principal $\mathcal W$-algebra
BV duality is established on the stated hook-type corridor; the
arbitrary-nilpotent extension asks for the corresponding BV map.
\item A Drinfeld-centre presentation of the closed sector asks for a
separate braided comparison with the categorical chiral center.
\item The Vol~III K3 row uses the Hall--Drinfeld/CoHA carrier for its
F5--F6 comparison.  The Drinfeld Yangian
$Y_\hbar(\mathfrak g_{\Delta_5})$ is a distinct carrier.  The constants
$K^{\kappa_{\mathrm{ch}}}(\mathrm{Mukai}(K3))=8$,
$\kappa_{\mathrm{BKM}}(\Delta_5)=5$,
$\kappa_{\mathrm{ch}}^{\mathrm{Heis}}(K3\times E)=3$, and
$\kappa_{\mathrm{cat}}(K3\times E)=0$ are distinct scalar invariants.
\end{enumerate}
